\documentclass[11pt]{article}

\usepackage[T1]{fontenc}
\usepackage[utf8]{inputenc}
\usepackage{amsmath,amssymb,mathtools}
\usepackage{amsthm}
\usepackage{array}
\newtheorem{lemma}{Lemma}
\newtheorem{proposition}{Proposition}
\usepackage{geometry}
\usepackage{enumitem}
\usepackage[expansion=false]{microtype}
\usepackage{natbib}
\usepackage[hidelinks]{hyperref}
\usepackage{draftwatermark}
\SetWatermarkText{In-Progress Draft}
\SetWatermarkScale{0.8}
\SetWatermarkColor[gray]{0.85}
\SetWatermarkAngle{53}

\geometry{margin=1in}
\setlist{nosep}

\title{The Future Economics of LLMs: From Hyperscaler GPUs to Consumer Usage}
\author{David Manheim}
\date{July 2026}

\begin{document}

\maketitle

\begin{abstract}
The model presented in this paper is straightforward, and novel only inasmuch as it applies standard market dynamics to the LLM market, and combines several parts of the market dynamics picture. Purchasers differ in how much they value model capability; models differ in economically relevant quality per unit of inference compute; hyperscalers have fixed short-run capacity; and model laboratories contract with those hardware owners for access to that capacity. We study the near-term dynamics of supply and demand which occur once purchasers want to buy more LLM usage than is available at the cost of production, so that prices can rise to reflect demand.

The central assumption is that LLM providers will increase prices to allow them to match demand, and thus scarce compute flows toward uses with the highest willingness to pay, which will effectively mean most LLM usage is for applications with the highest benefit to users per constrained FLOP. A laboratory's capability advantage matters only insofar as it raises this marginal net revenue.

There is also heterogeneity between lab capabilities; the persistent laboratory advantage is represented as a multiplicative productivity advantage rather than a fixed additive increment, which would vanish relative to an exponentially advancing frontier, and we account for heterogeneity across laboratory capabilities. The model shows that winner-take-all requires a demanding condition, which is that the leader's worst funded use of compute must be more valuable than every follower's best use. The paper then identifies several mechanisms that cause this condition to fail, producing coexistence instead. We distinguish the integrated-industry allocation benchmark from a plausible oligopolistic downstream market in which purchaser-facing prices contain three components: physical production cost, hyperscaler scarcity rent, and a model-provider markup. Finally, we extend the model to lumpy model development, temporary leadership, endogenous GPU investment, and alternative long-run technological regimes. Under a common asymptote, model-specific rents tend to disappear, though oligopoly markups need not; firm-specific asymptotes preserve durable differentiation; continuing improvement with a multiplicative advantage produces repeated episodes of temporary leadership in which laboratory rents rise after breakthroughs and hyperscaler bargaining power rises as competitors catch up; and continuing improvement with only an additive advantage combines unbounded quality growth with converging relative positions, so laboratory rents erode as a share of a growing market even though quality never stops improving.
\end{abstract}

\section{Introduction}
\label{sec:introduction}

Much discussion of AI market dynamics has simplified into pure competition by model providers, ignoring the fact that market for large-language-model inference has two economically distinct layers. Model laboratories own or control systems that differ in capability, efficiency, and commercial attractiveness. Hyperscalers own much of the hardware needed to operate those systems at scale. In the short run of multiple years, the stock of suitable inference hardware is fixed, while demand for valuable model outputs seems poised to become much larger than the amount that can be supplied at prices near physical production cost.

This setting raises four related questions. First, how is scarce inference compute allocated among models that differ in capability and cost? Second, how are the resulting rents divided between model laboratories and hardware owners? Third, does a persistent model advantage imply that one laboratory serves the entire market, or can trailing laboratories remain viable? Fourth, how do model development, GPU investment, and temporary frontier leadership change the static allocation over time?

The model below gives the obvious microeconomic answer, which serves as a baseline; compute is allocated according to marginal net benefit and willingness to pay\footnote{This does not require the usual assumption that compute flows only to the applications with the highest marginal revenue: some valuable consumer uses are still served, but the allocation implies that many consumer applications will be filled by relatively cheap models. We use marginal net revenue as a stand-in for marginal net benefit throughout.} per constrained FLOP\footnote{We do not explore labor market dynamics, but will note that as prices rise, companies will need to make much more conscious decisions about their LLM usage and spending; in some cases it's plausible that human programmers will temporarily remain cheaper than LLMs due to prices. This depends heavily on our later discussion about the future trajectory of model capabilities, but again, will not be explored in detail.}. A laboratory's capability advantage matters only insofar as it shifts this marginal-value schedule upward. Fixed hyperscaler capacity generates scarcity rents, but the leader need not receive all capacity because marginal demand slopes downward and competing models may generate greater value per FLOP for some tasks or purchaser segments.

The integrated-industry benchmark is useful but incomplete. When laboratories have downstream market power, purchaser-facing prices include a model-provider markup in addition to physical cost and the price of compute. Oligopoly can therefore leave consumer prices high even when GPU capacity is no longer scarce (Appendix~\ref{app:slack-capacity} works out that slack-capacity regime in full). Conversely, when capacity is tightly binding, oligopoly may change the division of rents between laboratories and hyperscalers more than it changes total output.

The dynamic extension treats training investment as competing with current inference for compute, makes model releases lumpy, and allows hyperscalers to expand future capacity through capital expenditure. A successful model release shifts a laboratory's marginal-value schedule upward and can generate temporary winner-take-most outcomes until a competitor releases a substitute. The duration and value of leadership depend on the rate of frontier progress, the persistence of firm-specific productivity differences, and the strategic behavior of oligopolistic model providers.

Throughout, we treat model laboratories and hyperscalers as distinct firms that transact at a price for compute. This is a simplification, and arguably the model's largest departure from the current industry: several of the largest hyperscalers own or hold major stakes in frontier laboratories, and leading laboratories increasingly build or contract for dedicated data-center capacity of their own. We keep the two layers separate because doing so makes the division of rents between them explicit. Most of the allocation results are, in any case, unchanged by integration: an integrated firm still faces the internal opportunity cost of a FLOP, so the allocation rule and the shadow price of capacity survive, and it is mainly the way the resulting rents are booked between the two layers that changes. Section~\ref{sec:vertical-integration} makes this precise.

\section{Literature Review}
\label{sec:literature-review}

The model is novel, if at all, in combining several standard mechanisms in the context of the LLM inference market. Its treatment of heterogeneous purchasers and differences in model quality follows the literature on vertical product differentiation.

\citet{mussaRosen1978} study a monopolist offering products of different qualities to consumers who differ in their willingness to pay for quality. \citet{shakedSutton1982} extend this logic to oligopolistic competition and show how vertical differentiation can soften price competition. The utility specification used here, in which purchaser $i$ values an output of quality $Q_m$ at $\theta_i \cdot Q_m$, is a reduced-form version of this standard approach. \citet{sutton1991} adds the entry margin: when the sunk cost of entry is endogenous, because firms can escalate quality investment, growing markets need not fragment and concentration stays bounded below. Section~\ref{sec:natural-oligopoly} applies this free-entry logic with training compute as the sunk cost.

The distinction between the integrated-industry benchmark and downstream oligopoly follows the broader literature on differentiated-product competition. \citet{dixit1979quality} analyzes quality and quantity choices under oligopoly, while \citet{singhVives1984} compare price and quantity competition in a differentiated duopoly. This literature motivates treating model-provider markups as distinct from physical production costs and upstream compute rents. It also shows why convergence in technical quality does not by itself imply marginal-cost pricing when a finite number of differentiated firms continue to compete strategically.

A second literature studies competition under capacity constraints. \citet{krepsScheinkman1983} show that capacity precommitment followed by price competition can yield the outcome of a Cournot quantity game. More directly related to the present model, \citet{boccardWauthy2010} combine vertical differentiation with capacity precommitment and analyze how product quality and limited capacity jointly affect price competition. These models establish that fixed capacity can sustain prices above marginal production cost and affect the allocation of demand across differentiated suppliers. The present model adapts this logic to a setting in which the constrained input is inference compute and models differ in the economic value they generate per FLOP.

The division of rents between laboratories and hyperscalers is related to the literature on vertical markets and bilateral bargaining. \citet{hornWolinsky1988} combine bargaining over input prices with downstream duopoly and show that the division of surplus depends on upstream and downstream market structure, product substitutability, and the parties' outside options. This provides a standard basis for separating the wholesale price of compute from the downstream model-provider markup, and for allowing both oligopoly among laboratories and market power among hyperscalers.

The dynamic behavior of trailing laboratories is related to models of learning-by-doing and dynamic pricing. \citet{cabralRiordan1994} show that current production and pricing can affect future competitive advantage through a learning curve. A firm may therefore price aggressively in the present because greater current output improves its future costs or market position. The analogous mechanism here is that serving inference demand may preserve customer relationships, generate deployment knowledge or data, and increase a laboratory's probability of producing a competitive future model. The strength of this mechanism depends on whether inference deployment actually contributes to future model capability.

The model-development race is also related to quality-ladder and creative-destruction models. \citet{grossmanHelpman1991} model repeated stochastic improvements that move products up a quality ladder, while \citet{aghionHowitt1992} study innovation driven by the temporary rents created when a new technology displaces an incumbent technology. These mechanisms correspond closely to the temporary leadership considered here: a successful model release raises a laboratory's productivity and market share until a competitor releases a sufficiently strong substitute. More directly, the investment-hazard structure used in Section~\ref{sec:model-race}, in which research spending raises the probability of a discrete release, follows the patent-race literature of \citet{loury1979}, \citet{leeWilde1980}, and \citet{reinganum1983}, while the leader's incentive to invest partly in order to deter followers is the preemption motive of \citet{gilbertNewbery1982}, and the observation that a follower gains more from a given release than the leader does is the replacement effect of \citet{arrow1962}. \citet{besankoDoraszelski2004} provide a complementary dynamic model of capacity accumulation in which lumpy investment, depreciation, and the mode of product-market competition can generate persistent asymmetries in firm size. Their analysis motivates the treatment of hyperscaler capacity as a slowly adjusting state variable rather than an instantaneously supplied input.

The subscription extension follows the literature on two-part tariffs. \citet{oi1971} considers a fixed access fee combined with a marginal usage charge, and \citet{schmalensee1981} analyzes how consumer heterogeneity affects optimal two-part pricing. The same structure applies when a user pays a subscription fee for the option to purchase future model outputs at a discounted marginal price. The extension does not require uncertainty: task-specific differences in model value can make several subscriptions worthwhile under complete information. When quality is costly to determine before use, the extension also relates to \citet{nelson1970}, who distinguishes goods whose relevant quality is learned through search or experience.

The comparison among technological regimes is related both to the economics of research productivity and to empirical neural scaling laws. \citet{kaplanEtAl2020} document power-law relationships between language-model loss and model size, data, and training compute. \citet{hoffmannEtAl2022} refine the compute-allocation problem by jointly scaling model size and training data. These results support modeling frontier progress as a systematic function of effective training compute over the observed range, but do not determine whether the relationship continues indefinitely or approaches a binding economic or technical asymptote. \citet{bloomEtAl2020} document declining research productivity in several sectors despite rising research effort, motivating the comparison between continuing improvement and cases in which additional performance gains require rapidly increasing investment.

A recent\footnote{or at least, recent as the term is used in economics; the AI literature moves much more quickly.} AI-specific literature studies model pricing, product design, and market structure. \citet{mahmood2024} analyzes pricing competition between generative-AI providers whose models differ in their cost-effectiveness across tasks. Users choose between models based on both price and the number of interactions required to obtain a satisfactory output. This directly addresses competition among differently performing models, but does not impose a common fixed supply of upstream inference compute.

\citet{bergemannBonattiSmolin2026} study optimal pricing and product design for LLM services when users differ in their valuations across tasks. Their framework models token-processing costs, heterogeneous users, and the allocation of inference resources, but focuses on product design by a single provider rather than competition among laboratories for a fixed stock of hyperscaler capacity.

At the industry level, \citet{korinekVipra2025} analyze how scaling, high fixed costs, access to compute and other scarce inputs, vertical integration, and feedback effects may produce concentration in the market for foundation models. They emphasize possible tipping near the technological frontier while allowing stronger competition among models behind it. \citet{demirerEtAl2025} provide complementary empirical evidence on model supply, prices, performance, demand, and differentiation. \citet{fradkin2025demand} documents rapid adoption following model releases, variation in whether new models substitute for existing demand or expand the market, and substantial multi-homing among applications. Together, these papers motivate treating model releases as discrete shocks to market shares while allowing several models to retain demand at the same time.

The present model combines heterogeneous willingness to pay for model quality, differences in quality-adjusted output per FLOP, fixed aggregate hyperscaler capacity, allocation according to marginal net revenue per FLOP, persistent laboratory-specific productivity differences, rent division between laboratories and hardware owners, and the possibility that trailing laboratories operate at a current loss to preserve future option value. The extensions add oligopolistic downstream pricing, two-part subscription contracts, lumpy quality-ladder innovation, endogenous hyperscaler capacity investment, and alternative long-run productivity regimes.

\section{Purchasers, Models, and Compute Requirements}
\label{sec:basic-setup}

There is a continuum of purchasers indexed by $i$, with total mass $N$. Purchaser $i$ has willingness-to-pay parameter $\theta_i>0$, drawn from a distribution $F$.

There are $M$ models. Model $m$ produces one standardized output of economically relevant quality $Q_m>0$ using an expected $a_m>0$ FLOPs under a specified workload and inference policy. An output may be interpreted as a fixed number of generated tokens, a completed query, or another standardized inference service.

Define quality-adjusted productivity per FLOP as

\begin{equation}q_m\equiv\frac{Q_m}{a_m}.\end{equation}

The choice of standardized output unit is a normalization: Appendix~\ref{app:extensions} shows that the resulting allocation is invariant to the choice of output unit, so nothing below depends on how a standardized output is defined.

Let $P_m$ denote the purchaser-facing price of one standardized output, and define the corresponding price per FLOP as

\begin{equation}p_m\equiv\frac{P_m}{a_m}.\end{equation}

Purchaser utility is determined by total output quality and total output price:

\begin{equation}u_{im}=\theta_i\cdot Q_m-P_m=a_m\cdot\left(\theta_i\cdot q_m-p_m\right).\end{equation}

Each purchaser chooses the model providing the highest non-negative utility:

\begin{equation}m_i^*=\arg\max_{m\in\{1,\ldots,M\}}\left\{0,\theta_i\cdot Q_m-P_m\right\}.\end{equation}

The normalized quantities $q_m$ and $p_m$ are useful for studying the allocation of compute, but purchasers compare total service bundles. In particular, when $a_m$ differs across models, maximizing $a_m\cdot(\theta_i\cdot q_m-p_m)$ is not equivalent to maximizing $\theta_i\cdot q_m-p_m$.

Let $\mathbf y=(y_1,\ldots,y_M)$ denote quantities of model outputs and let $\mathbf Q=(Q_1,\ldots,Q_M)$ denote their qualities. Define $R(\mathbf y;\mathbf Q)$ as the maximum total downstream revenue obtainable from selling output vector $\mathbf y$. We assume that $R$ is differentiable, increasing, and concave over the relevant capacity-constrained region. Concavity reflects the need to serve progressively lower-value purchasers as output expands.

The revenue function is an integrated-industry object: it represents coordinated downstream pricing or an equivalent implementation through contracts. A decentralized equilibrium among competing laboratories need not maximize $R$. The benchmark remains useful because it identifies the joint-profit-maximizing allocation of a scarce input independently of how the resulting rents are divided. This is a joint-profit benchmark rather than a welfare benchmark: because it equalizes marginal net revenue rather than marginal social value per FLOP, it coincides with the welfare-maximizing (efficient) allocation only under perfect price discrimination or price-taking behavior.

\begin{table}[htbp]
\centering
\caption{Main recurring symbols.}
\label{tab:notation}
\begin{tabular}{l>{\raggedright\arraybackslash}p{4.4in}}
\hline
Symbol & Meaning \\
\hline
$\theta_i$ & Purchaser $i$'s willingness-to-pay parameter for quality \\
$Q_m$ & Economically relevant quality of one standardized output of model $m$ \\
$a_m$ & FLOPs required per standardized output of model $m$ \\
$q_m\equiv Q_m/a_m$ & Quality-adjusted productivity per FLOP \\
$P_m$ & Purchaser-facing price of one standardized output \\
$p_m\equiv P_m/a_m$ & Purchaser-facing price per FLOP \\
$K$ & Aggregate short-run inference capacity (FLOPs per period) \\
$x_m$ & Compute (FLOPs) allocated to model $m$ \\
$y_m\equiv x_m/a_m$ & Output quantity of model $m$ \\
$c$ & Physical operating cost per FLOP \\
$b_m$ & Other variable cost per standardized output of model $m$ \\
$\lambda$ & Shadow price of the capacity constraint \\
$v_m$ & Marginal net revenue per FLOP for model $m$ (competitive benchmark) \\
$\eta_\ell$ & Persistent laboratory-specific productivity-level multiplier \\
$\alpha_\ell$ & Laboratory advantage increment: log-productivity term with $\eta_\ell=e^{\alpha_\ell}$ in the multiplicative specification; additive level increment in the additive specification (Secs.~\ref{sec:persistent-advantage}, \ref{subsec:additive-frontier}) \\
$r^*$ & Market-clearing wholesale value of a FLOP \\
$\mu_\ell$ & Downstream oligopoly markup of laboratory $\ell$ \\
$\widehat v_\ell$ & Laboratory $\ell$'s private marginal value per FLOP under oligopoly \\
$\zeta$ & Laboratory Nash bargaining weight (Appendix~\ref{app:bargaining}) \\
$\sigma$ & Compute-share adjustment speed after a release (Sec.~\ref{sec:model-race}) \\
$T_{\mathrm{build}}$ & Construction lag for new hyperscaler capacity \\
$\psi_\ell$ & Laboratory research-efficiency (rate-of-approach) parameter (Sec.~\ref{sec:technology-regimes}) \\
$\chi_{\ell,k}$ & Laboratory- and task-specific comparative-advantage term (Sec.~\ref{sec:persistent-advantage}) \\
$\rho$ & Discount rate (Secs.~\ref{sec:model-race}, \ref{sec:natural-oligopoly}); distinct from the wholesale compute price $r,r^*$ \\
\hline
\multicolumn{2}{l}{\emph{Symbols reused for unrelated objects, disambiguated by context or a note at first use:}} \\
$A_t$ vs.\ $A$ & Log frontier level (Sec.~\ref{sec:persistent-advantage}) vs.\ linear-demand intercept (Prop.~\ref{prop:capacity-neutrality} and after) \\
$G_\ell(x)$ vs.\ $G_{\ell t}$ & Gross operating value (Sec.~\ref{sec:concentration}) vs.\ cumulative training compute (Sec.~\ref{sec:technology-regimes}); $\Gamma_\ell$ is used for the entry-cost inverse to avoid a third use \\
$F$ (distribution) vs.\ $F$ (follower) & Sec.~\ref{subsec:two-model-demand} \\
$r^*$ (competitive benchmark) vs.\ $r^*$ (oligopoly auction) & footnote at first oligopoly use, Sec.~\ref{sec:oligopoly} \\
$r$ (wholesale compute price) vs.\ $a_m^r$, $\lambda_r$, $c_r$ (resource index) & Appendix~\ref{app:extensions}; the resource-index use is always a super/subscript on a different letter \\
\hline
\end{tabular}
\end{table}

\section{Fixed Inference Capacity}
\label{sec:fixed-capacity}

Inference hardware is owned by hyperscalers indexed by $h$. Hyperscaler $h$ has fixed short-run capacity $K_h$, measured in FLOPs per period. Aggregate capacity is

\begin{equation}K=\sum_h K_h.\end{equation}

Let $x_{hm}$ denote the FLOPs supplied by hyperscaler $h$ to model $m$. Each hyperscaler faces

\begin{equation}\sum_{m=1}^{M}x_{hm}\leq K_h.\end{equation}

Total compute allocated to model $m$ is

\begin{equation}x_m=\sum_h x_{hm},\end{equation}

and model output is

\begin{equation}y_m=\frac{x_m}{a_m}.\end{equation}

The aggregate capacity constraint is therefore

\begin{equation}\sum_{m=1}^{M}a_m\cdot y_m\leq K.\end{equation}

Let $c$ denote the physical operating cost per FLOP, and let $b_m$ denote other variable costs per standardized output from model $m$.

The maintained scarcity assumption is not that demand exceeds capacity at every possible price. At sufficiently high prices, demand falls. Instead, we study the parameter region in which demand for the relevant class of inference services exceeds available capacity at prices close to physical production cost.

Define the near-cost price per output as

\begin{equation}P_m^0=b_m+a_m\cdot c,\end{equation}

and let $D_m(\mathbf P,\mathbf Q)$ denote demand for model $m$. The scarcity assumption is

\begin{equation}\sum_{m=1}^{M}a_m\cdot D_m(\mathbf P^0,\mathbf Q)>K.\end{equation}

Under this assumption, the market cannot serve all demand at prices close to marginal physical cost. Capacity binds in the relevant benchmark allocation, and equilibrium prices may contain a scarcity rent.

The opposite case is a genuine possibility rather than a degenerate one: if desired output at near-cost prices falls short of $K$, the wholesale price falls to physical cost and ordinary unconstrained competition reappears. The main text flags this regime wherever an individual result depends on capacity binding; Appendix~\ref{app:slack-capacity} collects the slack-capacity regime in one place and works out what the model predicts there.

\subsection{Demand and revenue from purchaser heterogeneity}
\label{subsec:two-model-demand}

So far $R(\mathbf y;\mathbf Q)$ and $D_m(\mathbf P,\mathbf Q)$ have been treated as abstract primitives. They can instead be derived directly from the purchaser utility $\theta_i\cdot Q_m-P_m$. The standard case is two vertically differentiated models and uniformly distributed willingness to pay, which we work out here because it grounds two objects used later: the concavity of $R$ imposed in Section~\ref{sec:basic-setup}, and the assumption $\partial v_m/\partial q_m>0$ introduced in Section~\ref{sec:persistent-advantage}. A notational note before proceeding: the uniform specification below is the case $F=U[0,\bar\theta]$ of the willingness-to-pay distribution $F$ introduced in Section~\ref{sec:basic-setup}, while the subscript $F$ used throughout this subsection and Appendix~\ref{app:two-model} (as in $D_F$, $Q_F$, $\widehat v_F$) denotes the follower model; the two uses of $F$ are unrelated, and context disambiguates them.

Let there be two models, a leader $L$ and a follower $F$, with $Q_L>Q_F>0$, and let $\theta_i\sim U[0,\bar\theta]$ over a purchaser mass $N$. Each purchaser buys the model giving the highest non-negative utility. Comparing $L$ with $F$ gives an upper indifference cutoff, and comparing $F$ with buying nothing gives a lower cutoff:

\begin{equation}\theta^{*}=\frac{P_L-P_F}{Q_L-Q_F},\qquad \theta^{0}=\frac{P_F}{Q_F}.\end{equation}

When $\theta^{0}<\theta^{*}<\bar\theta$, the $\theta$ line segments into three quality tiers: purchasers with $\theta_i\geq\theta^{*}$ buy $L$, those with $\theta_i\in[\theta^{0},\theta^{*})$ buy $F$, and those with $\theta_i<\theta^{0}$ buy nothing. Integrating the uniform density over each tier gives the demand system

\begin{equation}D_L(\mathbf P;\mathbf Q)=N\left[1-\frac{P_L-P_F}{(Q_L-Q_F)\,\bar\theta}\right],\qquad D_F(\mathbf P;\mathbf Q)=N\left[\frac{P_L-P_F}{(Q_L-Q_F)\,\bar\theta}-\frac{P_F}{Q_F\,\bar\theta}\right].\end{equation}

Because each purchaser buys one output, $y_m=D_m$; inverting the system gives the inverse demands

\begin{equation}P_L=\bar\theta\left[Q_L-\frac{Q_L\,y_L+Q_F\,y_F}{N}\right],\qquad P_F=\bar\theta\,Q_F\left[1-\frac{y_L+y_F}{N}\right],\end{equation}

and downstream revenue $R(\mathbf y;\mathbf Q)=P_L\,y_L+P_F\,y_F$ is the quadratic form

\begin{equation}R(\mathbf y;\mathbf Q)=\bar\theta\left(Q_L\,y_L+Q_F\,y_F\right)-\frac{\bar\theta}{N}\left(Q_L\,y_L^{2}+2Q_F\,y_L\,y_F+Q_F\,y_F^{2}\right).\end{equation}

\begin{lemma}\label{lem:two-model}
For the two-model system above with $Q_L>Q_F>0$ and $\theta_i\sim U[0,\bar\theta]$:
\begin{enumerate}
\item $R(\mathbf y;\mathbf Q)$ is strictly concave in $\mathbf y$ over the region in which all three tiers are non-empty; and
\item writing $Q_m=a_m\,q_m$ and holding output quantities fixed, marginal net revenue per FLOP satisfies $\partial v_m/\partial q_m>0$ whenever the served mass $y_L+y_F$ is below $N/2$, i.e., whenever capacity is scarce enough that only the higher-$\theta$ portion of purchasers is served.
\end{enumerate}
\end{lemma}

Part~(1) is the microfoundation for the concavity assumption on $R$: the Hessian of $R$ is $-2(\bar\theta/N)\left[\begin{smallmatrix}Q_L&Q_F\\ Q_F&Q_F\end{smallmatrix}\right]$, which is negative definite precisely because $Q_L>Q_F$. Part~(2) grounds the productivity assumption of Section~\ref{sec:persistent-advantage}: a more productive model delivers higher economically relevant quality per FLOP, which raises its marginal revenue as long as it is not already forced down into the low-value part of the market. Appendix~\ref{app:two-model} gives the calculations, and shows that the same demand system delivers the comparative statics used in Section~\ref{sec:technology-regimes}.

\section{Integrated-Industry Allocation Benchmark}
\label{sec:allocation}

Suppose downstream quantities and prices can be coordinated, or equivalently that contracts among model providers and hyperscalers implement the allocation that maximizes aggregate operating profit. Transfers between model providers and hyperscalers do not affect this joint-profit objective.

Aggregate operating profit is

\begin{equation}W(\mathbf y)=R(\mathbf y;\mathbf Q)-\sum_{m=1}^{M}b_m\cdot y_m-c\cdot\sum_{m=1}^{M}a_m\cdot y_m.\end{equation}

The benchmark allocation solves

\begin{equation}\max_{\mathbf y\geq0}\;R(\mathbf y;\mathbf Q)-\sum_{m=1}^{M}b_m\cdot y_m-c\cdot\sum_{m=1}^{M}a_m\cdot y_m\quad\text{subject to}\quad\sum_{m=1}^{M}a_m\cdot y_m\leq K.\end{equation}

Let $\lambda\geq0$ be the multiplier on the capacity constraint. For every model receiving a positive allocation, the first-order condition is

\begin{equation}\frac{\frac{\partial R}{\partial y_m}-b_m}{a_m}=c+\lambda.\end{equation}

Define marginal revenue and marginal net revenue per FLOP as

\begin{equation}MR_m\equiv\frac{\partial R}{\partial y_m},\qquad v_m\equiv\frac{MR_m-b_m}{a_m}.\end{equation}

The allocation condition becomes

\begin{equation}\boxed{v_m=c+\lambda}\end{equation}

for every active model, while inactive models satisfy

\begin{equation}v_m\leq c+\lambda.\end{equation}

Thus scarce compute is allocated according to marginal net revenue per FLOP.

When the total output price can be treated as locally fixed because an additional output can be sold without materially changing price,

\begin{equation}MR_m\simeq P_m=a_m\cdot p_m.\end{equation}

Hence

\begin{equation}v_m\simeq\frac{P_m-b_m}{a_m}=p_m-\frac{b_m}{a_m}.\end{equation}

This is the sense in which laboratories and hyperscalers maximize dollars per FLOP. Strictly speaking, the relevant object is marginal net revenue per constrained FLOP, not average revenue per FLOP.

A more capable model does not necessarily receive compute before a weaker model. It does so only when its additional purchaser value is large enough, relative to its inference and non-compute costs, to generate greater marginal net revenue per FLOP. A lower-compute model has no independent advantage merely because its standardized output unit is smaller. It has an economic advantage when it preserves more than a proportional share of purchaser value or otherwise raises value net of cost per constrained resource.

\section{Persistent Laboratory Productivity}
\label{sec:persistent-advantage}

Each model is owned by a laboratory $\ell(m)$. Because economically relevant model productivity may improve by orders of magnitude, we model common technological progress and persistent laboratory advantage multiplicatively.

Let productivity per FLOP evolve according to

\begin{equation}\log q_{mt}=A_t+\alpha_{\ell(m)}+\varepsilon_{mt},\end{equation}

or equivalently,

\begin{equation}q_{mt}=\bar q_t\cdot\eta_{\ell(m)}\cdot\xi_{mt},\qquad \bar q_t\equiv e^{A_t},\qquad \eta_\ell\equiv e^{\alpha_\ell},\qquad \xi_{mt}\equiv e^{\varepsilon_{mt}}.\end{equation}

Here, $\bar q_t$ is the common technological frontier in economically relevant quality per FLOP, $\eta_\ell>0$ is a persistent laboratory-specific productivity multiplier, and $\xi_{mt}>0$ is a temporary model-specific shock. Total quality remains

\begin{equation}Q_{mt}=a_{mt}\cdot q_{mt}.\end{equation}

Laboratory $L$ has a persistent productivity lead over laboratory $j$ when $\eta_L>\eta_j$, or equivalently $\alpha_L>\alpha_j$. In the absence of temporary shocks,

\begin{equation}\frac{q_{Lt}}{q_{jt}}=\frac{\eta_L}{\eta_j}=e^{\alpha_L-\alpha_j}.\end{equation}

The proportional advantage is therefore independent of the level of the common frontier.

This differs importantly from an additive level specification,

\begin{equation}q_{mt}=A_t+\alpha_{\ell(m)}+\varepsilon_{mt}.\end{equation}

Under the additive specification,

\begin{equation}\lim_{A_t\rightarrow\infty}\frac{A_t+\alpha_L}{A_t+\alpha_j}=1.\end{equation}

A fixed laboratory increment becomes asymptotically negligible relative to the common frontier. This limit is set aside here because it cannot represent \emph{persistent} advantage, but it is not discarded: Section~\ref{subsec:additive-frontier} returns to it as a long-run technological regime in its own right---unbounded improvement with converging relative positions. The multiplicative specification instead represents a persistent structural efficiency advantage in algorithms, training, post-training, data, research organization, or the ability to produce economically relevant performance with fewer inference resources.

Higher productivity per FLOP generally shifts the model's marginal-value schedule upward. We assume

\begin{equation}\frac{\partial v_m}{\partial q_m}>0.\end{equation}

Lemma~\ref{lem:two-model}(2) derives this inequality for the two-model demand system, under the scarcity condition that only the higher-$\theta$ portion of purchasers is served; we take it as a maintained assumption more generally. It follows that

\begin{equation}\eta_L\uparrow\;\Longrightarrow\;q_L\uparrow\;\Longrightarrow\;v_L\uparrow.\end{equation}

A persistent proportional productivity advantage therefore creates a persistent advantage in the market for scarce compute. Whether it produces a growing absolute rent also depends on downstream demand, capacity growth, and the elasticity of marginal revenue with respect to productivity.

A useful alternative interpretation is a persistent lead in frontier time. If $\bar q_t=\bar q_0\cdot e^{g\cdot t}$ and laboratory $\ell$ is effectively $\tau_\ell$ periods ahead, then

\begin{equation}q_{\ell t}=\bar q_{t+\tau_\ell}=\bar q_t\cdot e^{g\cdot\tau_\ell}.\end{equation}

Thus a fixed time lead corresponds to a productivity multiplier $\eta_\ell=e^{g\cdot\tau_\ell}$.

The decomposition above compresses a laboratory's productivity into the single multiplier $\eta_\ell$, which implicitly treats relative performance as uniform across tasks. That is a simplifying assumption, not a restriction of the underlying mechanism: the same multiplicative structure extends directly to task-indexed productivity. Let $k$ index tasks or purchaser segments and write

\begin{equation}\log q_{m,k,t}=A_t+\alpha_{\ell(m)}+\chi_{\ell(m),k}+\varepsilon_{m,k,t},\end{equation}

where $\chi_{\ell,k}$ is a laboratory- and task-specific comparative-advantage term (we use $\chi$ rather than the $\beta$ already carrying two other meanings in this paper---the hyperscaler discount factor $\beta^s$ of Section~\ref{sec:model-race} and the scaling-law curvature exponent $\beta$ of Section~\ref{sec:technology-regimes}---since $\chi_{\ell,k}$ is the newest and least entrenched of the three and is cheapest to rename). We fix a single, common, market-wide task weighting $w_k$ once---for concreteness, task $k$'s share of aggregate purchaser mass---and normalize $\sum_k w_k\,\chi_{\ell,k}=0$ under that common weighting for every laboratory $\ell$, leaving $\alpha_\ell$ as laboratory $\ell$'s productivity averaged against the same task mix used for every other laboratory. This common weighting is what makes $\alpha_L$ and $\alpha_j$ comparable across laboratories; weighting each laboratory by its own task mix instead would average different laboratories over different task distributions and break that comparison. (Because the normalization is imposed on log productivity, $\alpha_\ell$ is a weighted geometric-mean advantage rather than an arithmetic average of productivity levels; the two can diverge once the dispersion of $\chi_{\ell,k}$ across tasks is large, by Jensen's inequality.) A laboratory with $\alpha_L>\alpha_j$ can still trail on individual tasks where $\chi_{L,k}<\chi_{j,k}$ outweighs the average gap, which is exactly the source of the coexistence mechanisms taken up in Section~\ref{subsec:task-segmentation} and the task-specific productivity of Appendix~\ref{app:extensions}. We work with the scalar $\eta_\ell$ through Section~\ref{sec:concentration} as an expository simplification, not because it is adequate for winner-take-all: the aggregate condition of that section, $v_L(K)\geq\max_{\ell\neq L}v_\ell(0)$, is in general \emph{not} robust to task dispersion once $v_\ell$ is built up from a task-indexed schedule, because a follower's $v_\ell(0)$ is then governed by its \emph{best} task, $\max_k\chi_{\ell,k}$, not by the task-weighted average $\sum_k w_k\chi_{\ell,k}=0$ that pins down $\alpha_\ell$. A follower with low average productivity but one strong task can still clear $r^*$ on that task alone and break aggregate winner-take-all even though it would not on the task-weighted average---this is exactly the protected-niche mechanism that Proposition~\ref{prop:segment-wta} formalizes twenty lines below. Task-level detail is therefore not a refinement that matters only once concentration is already known to be incomplete; it is necessary whenever the \emph{identity} of the aggregate winner-take-all outcome, rather than only average productivity, is the object of interest.

\section{Concentration and Coexistence}
\label{sec:concentration}

To state the concentration result transparently, impose a reduced-form separability assumption at the laboratory level. Let $G_\ell(x_\ell)$ be the gross operating value generated by laboratory $\ell$ after it optimally allocates $x_\ell$ FLOPs among its own models. Assume

\begin{equation}G_\ell'(x_\ell)>0,\qquad G_\ell''(x_\ell)<0.\end{equation}

Define the laboratory's marginal net revenue per FLOP as

\begin{equation}v_\ell(x_\ell)\equiv G_\ell'(x_\ell).\end{equation}

The reduced form abstracts from strategic cross-effects in downstream pricing. Appendix~\ref{app:extensions} states the corresponding condition when laboratories' marginal values depend directly on rivals' allocations.

Let $r^*$ be the market-clearing wholesale value of a FLOP in the competitive-allocation benchmark. Every active laboratory satisfies

\begin{equation}v_\ell(x_\ell)=r^*,\end{equation}

while inactive laboratories satisfy

\begin{equation}v_\ell(0)\leq r^*.\end{equation}

Capacity clears:

\begin{equation}\sum_\ell x_\ell=K.\end{equation}

Let $L$ denote the laboratory with the highest initial marginal value per FLOP. It receives all available compute if and only if

\begin{equation}\boxed{v_L(K)\geq\max_{\ell\neq L}v_\ell(0).}\end{equation}

This condition says that even the leader's least valuable use of the final available FLOP is more valuable than every follower's best initial use.

If the condition fails, multiple laboratories receive compute. The model does not itself predict which outcome occurs: winner-take-all requires the demanding condition above---the leader's worst funded use of the final available FLOP must be more valuable than every follower's best initial use---and the mechanisms listed below cause that condition to fail, producing coexistence instead. Whether the condition holds is an empirical question about the curvature of laboratories' value schedules that the model does not restrict. Coexistence requires different marginal-value-per-FLOP schedules, which may arise because:

\begin{enumerate}
\item laboratories differ in quality-adjusted productivity per FLOP;
\item relative performance differs across tasks or purchaser segments (Section~\ref{subsec:task-segmentation} formalizes this mechanism, together with the frictions of the next);
\item models are horizontally differentiated by language, latency, privacy, specialization, reliability, integration, or contractual guarantees;
\item laboratories face different normalized non-compute costs or compute-contract terms; or
\item a laboratory offers several model tiers that serve different portions of demand.
\end{enumerate}

The final mechanism implies that concentration at the laboratory level may exceed concentration at the model level. A single laboratory may operate frontier, mid-tier, and lower-priced models. A smaller model contributes to the laboratory's compute demand only if it raises marginal net revenue per constrained FLOP for some portion of demand, not merely because one standardized output uses fewer FLOPs. All of this takes the set of laboratories as given: it asks how laboratories that already hold frontier models divide scarce compute. Section~\ref{sec:natural-oligopoly} treats the prior question of how many laboratories can recover the fixed cost of fielding such a model at all.

\subsection{Task-Indexed Quality and Segment-Level Dominance}
\label{subsec:task-segmentation}

The second and third mechanisms deserve more than a bullet, because the compression of commercially relevant performance into the scalar $Q_m$ (or $q_m$) is doing real work in the winner-take-all condition above. Index use cases or purchaser segments by $k$: coding, enterprise agents, low-latency consumer queries, a particular language. Let $Q_{mk}$ denote model $m$'s economically relevant quality on use case $k$, with $a_{mk}$ FLOPs per standardized output and $q_{mk}\equiv Q_{mk}/a_{mk}$, as in the task-specific productivity of Appendix~\ref{app:extensions}. Purchaser utility becomes

\begin{equation}u_{imk}=\theta_{ik}\cdot Q_{mk}-P_{mk}-s_{imk},\end{equation}

where $\theta_{ik}$ allows a purchaser's valuation of quality to differ across use cases and $s_{imk}\geq0$ is a purchaser- and provider-specific friction: an integration cost, latency, a privacy cost, a switching cost, or a contractual requirement. Setting a single use case and $s_{imk}=0$ recovers the baseline. The friction term is what absorbs the third mechanism: two models with identical $Q_{mk}$ can face different effective demand because purchasers bear different frictions in reaching them. One distinction within that informal list matters enough to flag now. Because $s_{imk}$ enters utility purchase by purchase, it is a \emph{recurring}, per-output cost: latency, a privacy cost, per-task contractual or compliance overhead---costs borne on every interaction with a non-preferred provider. A cost paid \emph{once} to establish a relationship with a provider---evaluating it, integrating it initially, negotiating terms---is a different primitive: it is the evaluation cost $E_{ij}$ of Appendix~\ref{app:subscriptions}, paid a single time to add a provider to a subscription portfolio, after which routing among subscribed providers there is frictionless. The informal list above mixes the two (``an integration cost\ldots a switching cost'' reads naturally as one-time), and they protect an incumbent very differently: a recurring cost protects it on every future purchase, while a one-time cost protects it only against purchasers who have not yet paid it, and stops protecting it entirely against any purchaser who multi-homes. The distinction is load-bearing wherever a friction magnitude does formal work, in particular in Proposition~\ref{prop:segmented-entry}. The setup is close in spirit to \citet{mahmood2024}, where models differ in cost-effectiveness across tasks, and to \citet{bergemannBonattiSmolin2026}, where users differ in valuations across tasks.

To carry this section's machinery over, we need the analogue of its reduced-form separability assumption, now at the laboratory--segment level. Assume laboratory $\ell$ generates gross operating value $\sum_k G_{\ell k}(x_{\ell k})$ from devoting $x_{\ell k}$ FLOPs to segment $k$, with $G_{\ell k}'>0$ and $G_{\ell k}''<0$ for each pair, and that compute is fungible across segments, so the aggregate constraint is $\sum_\ell\sum_k x_{\ell k}\leq K$. Write $v_{\ell k}(x)\equiv G_{\ell k}'(x)$. Separability is doing work at two different levels here, and only one of them is the cross-segment abstraction flagged when the laboratory-level version of this assumption was introduced. Across segments, it abstracts from linkages such as shared subscriptions, bundled enterprise contracts, or reputational spillovers from a strong position in one segment onto demand in another---the same restriction as before, one level down. Within a segment, however, it also assumes $G_{\ell k}$ depends only on laboratory $\ell$'s own allocation $x_{\ell k}$, and \emph{not} on a rival's allocation in the same segment. That is a materially stronger restriction, and it is not innocuous: it is exactly what the paper's own microfoundation violates in the case that determines whether a segment is winner-take-all. In the two-model demand system of Lemma~\ref{lem:two-model}, when two laboratories serve the same segment, $MR_L=\bar\theta[Q_L-(Q_Ly_L+Q_Fy_F)/N]$ depends directly on $y_F$ (Section~\ref{subsec:two-model-demand}), so $G_{Lk}$ cannot be written as a function of $x_{Lk}$ alone once a rival is active in segment $k$. The separability assumption of this proposition is therefore valid precisely for segments served by one active laboratory against passive or inactive rivals---the case the proposition is really about---and is not a description of two laboratories actively contesting a single segment's demand. Proposition~\ref{prop:segment-wta} below should be read with that scope; the nonseparable case is treated only where the general condition of Appendix~\ref{app:extensions} is invoked explicitly.

\begin{proposition}[Segment-level winner-take-all]
\label{prop:segment-wta}
Under the separability assumption above---restricted, per the discussion just given, to configurations in which at most one laboratory is active in any given segment, with rivals passive or inactive there---let $r^*$ be the wholesale value of a FLOP that clears the capacity market, $\sum_\ell\sum_k x_{\ell k}=K$. Then:
\begin{enumerate}
\item every active laboratory--segment pair satisfies $v_{\ell k}(x_{\ell k})=r^*$, and every inactive pair satisfies $v_{\ell k}(0)\leq r^*$;
\item segment $k$ is served by laboratory $L_k$ alone if and only if $v_{L_k,k}(0)>r^*$ \emph{and} $\max_{\ell\neq L_k}v_{\ell k}(0)\leq r^*$: the winner-take-all condition of this section applies segment by segment, with the market-wide clearing value $r^*$ in the role of the leader's marginal value of the final FLOP. The participation condition $v_{L_k,k}(0)>r^*$ is not redundant: without it, a segment that no laboratory serves at all---every pair inactive, including $L_k$---also satisfies $\max_{\ell\neq L_k}v_{\ell k}(0)\leq r^*$, so the condition on rivals alone is necessary but not sufficient for sole service by $L_k$; and
\item the identity of $L_k$ may differ across $k$. Following the aggregate condition of Section~\ref{sec:concentration} rather than re-invoking $r^*$---whose supporting value is not unique at a corner where one laboratory absorbs everything---one laboratory $L$ absorbs all of $K$ if and only if $v_L(K)\geq\max_k\max_{\ell\neq L}v_{\ell k}(0)$, where $v_L(K)\equiv G_L'(K)$ is the common marginal value laboratory $L$ attains when it optimally allocates $K$ across every segment it could serve, i.e., the multiplier on $\sum_k x_{Lk}=K$ in $\max_{\{x_{Lk}\}}\sum_k G_{Lk}(x_{Lk})$.
\end{enumerate}
\end{proposition}

The proof is the first-order condition of Section~\ref{sec:allocation} applied to the doubly indexed choice variables $x_{\ell k}$, with strict concavity of each $G_{\ell k}$ delivering uniqueness; parts~(2) and~(3) read off which pairs clear the participation threshold $r^*$, with part~(2) requiring in addition that $L_k$ itself clear it. Part~(2) is the boxed condition of this section restated within a segment: the segment leader's least valuable funded FLOP in segment $k$ earns exactly $r^*$, so sole service of the segment requires both that $L_k$ actually wants to serve it and that every rival's best use \emph{in that segment} be worth no more than that. Part~(3) mirrors the aggregate boxed condition of Section~\ref{sec:concentration} exactly, comparing $L$'s own marginal value of the last unit of $K$ directly against every rival's value at zero in every segment, rather than selecting a particular supporting $r^*$ at a degenerate corner where many prices support the same allocation.

The qualitative headline is part~(3): nothing forces $L_k$ to be constant across $k$. Laboratory $L$ may lead on coding and enterprise agents while laboratory $j$ leads on low-latency consumer queries or on a particular language, and both positions are consistent with the same clearing value $r^*$. Global winner-take-all now requires the demanding condition of this section to hold for the same laboratory in every segment at once, and it fails whenever any rival holds a protected niche---a segment where task-specific quality or rivals' frictions give it $v_{\ell k}(0)>r^*$. Four implications follow, of which the first two are direct consequences of Proposition~\ref{prop:segment-wta} and the last two are plausible readings of it rather than results.

\paragraph{Aggregate versus within-segment concentration.} Because global dominance requires one laboratory to win every segment while segment-level dominance requires winning only one, aggregate concentration is weakly lower than a one-quality model calibrated to the same aggregate value schedules would predict, while concentration \emph{within} an individual use case can simultaneously be extreme. Measured market structure therefore depends on the unit of observation: an aggregate concentration index can look moderate while nearly every individual use case is served by one or two providers. The testable form is that concentration measured within a use case should exceed concentration measured across the market as a whole.

\paragraph{Persistent multi-homing.} A purchaser whose task mix spans segments led by different laboratories rationally pays for access to more than one provider. This is exactly the mechanism formalized in Appendix~\ref{app:subscriptions}, whose incremental-value condition requires only that a provider be sufficiently better or cheaper on \emph{some} tasks; task-indexed quality supplies the reason it is better, and frictions supply the reason routing does not instantly consolidate---with the division of labor just flagged: the one-time evaluation cost $E_{ij}$ governs which providers enter a purchaser's portfolio at all, while the recurring component of $s_{imk}$ governs how sticky routing remains among providers already in it. The multi-homing that \citet{fradkin2025demand} documents is, on this reading, an equilibrium feature of segmented leadership rather than a transitional artifact of immature tooling.

\paragraph{Less abrupt displacement.} In the one-quality model a strong release moves a single boxed condition, so market shares can jump discretely (Section~\ref{sec:model-race}). Here a release improves a laboratory's position segment by segment, and displacement occurs only in the segments where the per-segment threshold is crossed, so aggregate share should move in smaller steps. We flag this as plausible rather than proven: it depends on releases improving quality non-uniformly across segments, which is realistic but not something the model restricts, and a release that dominates uniformly still displaces everywhere at once. The mechanism is complementary to, and observationally hard to separate from, the switching-and-integration friction $\sigma$ of Section~\ref{sec:model-race}.

\paragraph{A portfolio interpretation of laboratory advantage.} A laboratory's competitive position is then the profile of its segment-level positions rather than the scalar $\eta_\ell$ of Section~\ref{sec:persistent-advantage}, its flow rents sum across segments, and the temporary-leadership value of Section~\ref{sec:model-race} becomes a sum of segment-level terms, each with its own displacement hazard. This is an interpretation of the model rather than a theorem, but it matters for reading the dynamics: a laboratory can lose the frontier in one segment while its aggregate position barely moves.

Differentiation of this kind also has a pricing consequence, which anticipates Section~\ref{sec:oligopoly}. Under downstream oligopoly the Lerner condition there applies segment by segment: with segment-level demand and own-price elasticity $\epsilon_{\ell\ell,k}$, the markup $\mu_{\ell k}$ is segment-specific, and it is largest in the segments where the laboratory's quality lead or its rivals' friction disadvantage is largest, because that is where competing models are the poorest substitutes and $\epsilon_{\ell\ell,k}$ is lowest. This is not, however, an instance of Proposition~\ref{prop:segment-wta} itself: the two-model demand system of Lemma~\ref{lem:two-model}, applied within a single segment with two laboratories actively contesting it, is precisely the nonseparable case that Proposition~\ref{prop:segment-wta}'s separability assumption sets aside (see the discussion opening this subsection). What the two-model system \emph{does} deliver exactly, directly and without appeal to Proposition~\ref{prop:segment-wta}, is Proposition~\ref{prop:breakthrough}'s markup comparative static in the quality gap, which is proven for that demand system on its own terms; in a general demand system the segment-specific-markup claim above is directional rather than guaranteed, for the reasons discussed before Proposition~\ref{prop:breakthrough}. The price decomposition of Section~\ref{sec:oligopoly} then holds segment by segment: the scarcity rent component reflects the common clearing value of fungible compute, while the markup component varies across use cases, so a single provider's price relative to quality should be highest precisely where its lead is largest.

Finally, this demand-side force and the supply side push in opposite directions: task-indexed demand supports more viable frontier providers, while vertical integration and dependence on particular clouds or accelerator families can reduce the number of laboratories able to field frontier capacity at all; see Section~\ref{sec:vertical-integration} for that upstream side. Segmentation also matters upstream of the rent division studied here: Section~\ref{subsec:segmented-entry} shows that defensible segments change not just how rents divide among existing laboratories but how many laboratories can recover the fixed cost of training in the first place.

\section{Oligopolistic Model Provision}
\label{sec:oligopoly}

The preceding allocation benchmark suppresses downstream market power. In an oligopoly, each laboratory chooses its own output or price while taking rivals' decisions as given. This introduces a model-provider markup in addition to physical production cost and the hyperscaler's compute charge.

Let laboratory $\ell$ face residual inverse demand

\begin{equation}P_\ell=P_\ell(y_\ell,\mathbf y_{-\ell};\mathbf q).\end{equation}

If laboratory $\ell$ pays wholesale compute price $r_\ell$ per FLOP, its marginal cost per standardized output is

\begin{equation}\kappa_\ell=b_\ell+a_\ell\cdot r_\ell.\end{equation}

Under Cournot quantity competition, operating profit is

\begin{equation}\pi_\ell=\left[P_\ell(y_\ell,\mathbf y_{-\ell};\mathbf q)-b_\ell-a_\ell\cdot r_\ell\right]\cdot y_\ell.\end{equation}

The interior first-order condition is

\begin{equation}P_\ell+y_\ell\cdot\frac{\partial P_\ell}{\partial y_\ell}=b_\ell+a_\ell\cdot r_\ell.\end{equation}

Define the laboratory's downstream oligopoly markup as

\begin{equation}\mu_\ell\equiv-y_\ell\cdot\frac{\partial P_\ell}{\partial y_\ell}>0.\end{equation}

The purchaser-facing price decomposes as

\begin{equation}\boxed{P_\ell=\underbrace{b_\ell+a_\ell\cdot c}_{\text{physical production cost}}+\underbrace{a_\ell\cdot(r_\ell-c)}_{\text{hyperscaler scarcity rent}}+\underbrace{\mu_\ell}_{\text{model-provider markup}}.}\end{equation}

This decomposition separates the two layers of market power. The hyperscaler earns a margin when compute is scarce or when hardware supply is concentrated. This margin is a peak-load scarcity rent of the kind first analyzed for capacity-constrained utilities by \citet{boiteux1960} and studied in the modern literature on scarcity pricing of fixed capacity \citep{crewKleindorfer1976}: with short-run capacity fixed and demand pressing against it, price rises above physical operating cost, and the resulting rent is the market's signal for capacity investment. The laboratory earns a downstream markup when its model faces a downward-sloping residual demand curve.

The laboratory's willingness to pay for an additional FLOP is no longer the average purchaser-facing margin per FLOP. It is the increase in equilibrium profit generated by additional compute:

\begin{equation}\label{eq:vhat-def}\widehat v_\ell(x_\ell,\mathbf x_{-\ell})\equiv\frac{\partial\pi_\ell^*}{\partial x_\ell}.\end{equation}

Two value concepts are used from here on. Plain $v$ (with subscript $m$ or $\ell$) denotes the integrated-benchmark, competitive marginal net value per FLOP of Sections~\ref{sec:allocation}--\ref{sec:concentration}, while the hatted $\widehat v_\ell$ denotes a laboratory's private marginal value per FLOP under oligopoly, used in this section and afterward. Because additional oligopoly output reduces the price earned on the laboratory's inframarginal units, $\widehat v_\ell\leq v_\ell$ in general.

Because additional output may reduce the price earned on the laboratory's existing output, generally

\begin{equation}\widehat v_\ell<\frac{P_\ell-b_\ell}{a_\ell}.\end{equation}

Under a competitive upstream auction for fixed capacity, the wholesale compute price $r^*$\footnote{This is the same symbol $r^*$ introduced in Section~\ref{sec:concentration} for the competitive-benchmark clearing price. Here it denotes the analogous market-clearing compute price under the oligopoly/auction setting; the two are distinct equilibrium objects that share one symbol.} is approximately determined by the best losing laboratory's private marginal value:

\begin{equation}r^*\simeq\max\{c,\widehat v_F\}.\end{equation}

The marginal surplus then decomposes as

\begin{equation}\widehat v_L-c=\left(\widehat v_L-\widehat v_F\right)+\left(\widehat v_F-c\right).\end{equation}

The first term is differential rent retained by the leading laboratory; the second is the hyperscaler's scarcity rent under the competitive upstream benchmark. Appendix~\ref{app:bargaining} relaxes this division by allowing bilateral bargaining and oligopsony. There, with a laboratory Nash bargaining weight $\zeta\in[0,1]$, the negotiated wholesale price splits the match-specific surplus $\widehat v_\ell-\bar v_h$ so that the laboratory retains a share $\zeta$ and the hyperscaler the remaining $1-\zeta$; the competitive-auction benchmark above is the special case in which the hyperscaler's outside option approaches the best losing bidder's value. That appendix bargains over a single marginal block in isolation, following \citet{hornWolinsky1988} for motivation; when several laboratory--hyperscaler pairs bargain simultaneously, the appropriate solution concept is Nash-in-Nash, in which each pair's price is the Nash bargain taken as given the agreements reached by every other pair \citep{collardWexlerEtAl2019}.

\subsection{Binding and slack capacity}

Oligopoly has different effects depending on whether aggregate compute remains scarce after laboratories restrict output. Let $y_\ell^{\mathrm{olig}}(c)$ denote laboratory $\ell$'s oligopoly output when compute is available at physical cost. Capacity remains binding if

\begin{equation}\sum_\ell a_\ell\cdot y_\ell^{\mathrm{olig}}(c)>K.\end{equation}

When this condition holds, all capacity is used. Downstream oligopoly may have a limited effect on total output because total quantity is physically constrained, but it changes the division of rents: laboratories bid less aggressively for marginal compute than competitive output providers would and retain larger downstream markups.

Capacity is slack if

\begin{equation}\sum_\ell a_\ell\cdot y_\ell^{\mathrm{olig}}(c)\leq K.\end{equation}

Then the wholesale compute price falls to physical cost, $r^*=c$, and hyperscaler scarcity rents disappear. Consumer prices may nevertheless remain above marginal cost because $\mu_\ell>0$. Additional GPU construction then has little direct effect on purchaser-facing prices unless it changes downstream entry, quality, or strategic conduct. Appendix~\ref{app:slack-capacity} works this regime out in full---the slack condition itself, the reappearing dependence of price on the number of laboratories, free entry, investment incentives near the boundary, and what remains of niche protection.

The binding case has a sharper implication, which the symmetric Cournot example of Appendix~\ref{app:cournot} makes exact. As with every proposition in this section, laboratories are price takers in the upstream compute market: they treat $r^*$ as parametric and it adjusts to clear $\sum_\ell a_\ell y_\ell=K$. Appendix~\ref{app:bargaining} relaxes this to bilateral Nash bargaining with laboratory weight $\zeta\in[0,1]$; under bargaining, $r^*$ in the results below is replaced by the negotiated price $r_{\ell h}=\zeta\bar v_h+(1-\zeta)\widehat v_\ell$ of that appendix, laboratories with bargaining power retain part of the scarcity rent that competitive pricing here transfers to hyperscalers, and the rent-transfer direction in part of the proposition below survives (raising $n$ still moves rent toward hyperscalers at any fixed $\zeta$) but its magnitude is smaller the larger is $\zeta$.

\begin{proposition}[Capacity-binding neutrality of the number of laboratories]
\label{prop:capacity-neutrality}
Consider $n$ symmetric laboratories selling a homogeneous inference service under Cournot competition with linear inverse demand $P(Y)=A-B\cdot Y$\footnote{This $A$ is the demand intercept of the linear inverse-demand system used in this proposition, Appendix~\ref{app:cournot}, and Propositions~\ref{prop:free-entry}--\ref{prop:segmented-entry}; it is unrelated to the log frontier level $A_t$ of Section~\ref{sec:persistent-advantage}. The two never appear in the same expression, and context (a subscript $t$, or its absence) disambiguates them throughout.}, where each standardized output requires $a$ FLOPs. If aggregate capacity binds, so that $a\cdot Y=K$, the purchaser-facing price is
\begin{equation}P^*=A-\frac{B\cdot K}{a},\end{equation}
which depends on aggregate capacity $K$ but not on the number of laboratories $n$. Raising $n$ leaves total output and the consumer price unchanged and instead reallocates rent: each laboratory's markup $B\cdot K/(a\cdot n)$ falls, while the wholesale compute price $r^*$ rises, transferring fixed-capacity rent from laboratories to hyperscalers. See Appendix~\ref{app:cournot} for the derivation.
\end{proposition}

The economic content is a clean separation of two policy levers. As long as capacity binds, building GPUs lowers purchaser prices, because added capacity relaxes the constraint that pins down output; adding laboratories does not, because total output is fixed by capacity and further competition only redistributes the resulting rent. Additional laboratories move the consumer price only once entry pushes total desired output below capacity, at which point capacity becomes slack, $r^*=c$, and the ordinary Cournot dependence on the number of firms reappears (derived in Appendix~\ref{app:slack-capacity}). The proposition is deliberately stated for the symmetric, homogeneous-good, linear-demand case: it isolates the mechanism rather than claiming that real laboratories are identical, and heterogeneity or differentiation reintroduces a role for the composition of providers even under binding capacity.

\subsection{Differentiated-price competition}
\label{subsec:price-competition}

This subsection and its use in Section~\ref{sec:technology-regimes} and Proposition~\ref{prop:breakthrough} are a scope departure from the rest of the paper, and worth flagging before the algebra. The paper's maintained regime, from Section~\ref{sec:fixed-capacity} onward, is that aggregate compute capacity binds. Unconstrained price competition under binding capacity is the classical Edgeworth problem: with demand pressing against a fixed physical limit, pure-strategy price equilibria generically fail to exist, which is exactly why \citet{krepsScheinkman1983}---cited in Section~\ref{sec:literature-review}---study capacity precommitment followed by price competition and recover Cournot outcomes instead. The Bertrand markup and elasticity analysis that follows, and that reappears in part~1 of Proposition~\ref{prop:breakthrough} (Appendix~\ref{app:two-model}), should therefore be read as a slack-capacity, or locally compute-inclusive-cost, analysis: it describes how laboratories compete on price when compute is obtainable at the prevailing compute-inclusive cost $\kappa_\ell$, not as a claim that unconstrained Bertrand describes the binding-capacity regime where the paper's main scarcity story lives. Where the paper needs a genuinely capacity-constrained result---the differential compute rent of Proposition~\ref{prop:breakthrough}, part~2---it uses the Kreps--Scheinkman capacity-then-price reduced form (a quantity game in the compute endowments $x_\ell$) rather than unconstrained Bertrand; Appendix~\ref{app:two-model} gives that derivation directly, so this caveat is a scope statement about which regime each result describes rather than an unresolved gap. Appendix~\ref{app:slack-capacity} treats the slack regime systematically rather than leaving it as a scope caveat.

If laboratories choose prices rather than quantities, let demand for laboratory $\ell$ be $D_\ell(\mathbf P;\mathbf q)$. Its operating profit is

\begin{equation}\pi_\ell=\left(P_\ell-\kappa_\ell\right)\cdot D_\ell(\mathbf P;\mathbf q).\end{equation}

The interior price-setting condition is

\begin{equation}D_\ell+\left(P_\ell-\kappa_\ell\right)\cdot\frac{\partial D_\ell}{\partial P_\ell}=0.\end{equation}

If $\epsilon_{\ell\ell}$ is the absolute own-price elasticity, the Lerner condition is

\begin{equation}\frac{P_\ell-\kappa_\ell}{P_\ell}=\frac{1}{\epsilon_{\ell\ell}}.\end{equation}

A frontier leader can charge a large markup when competing models are poor substitutes and its own-price elasticity is low. A strong competitor release increases substitution, raises the leader's elasticity, compresses its markup, and simultaneously improves the hyperscaler's outside option in the upstream compute market.

\section{Prices and the Division of Rents}
\label{sec:prices-rents}

The competitive-allocation benchmark remains useful for separating general scarcity rent from differential laboratory productivity, but the oligopoly section adds a third component: downstream markup. At the purchaser-facing level,

\begin{equation}P_\ell-\left(b_\ell+a_\ell\cdot c\right)=a_\ell\cdot(r_\ell-c)+\mu_\ell.\end{equation}

At the margin, under a competitive upstream auction,

\begin{equation}\widehat v_L-c=\underbrace{\widehat v_L-\widehat v_F}_{\text{laboratory differential rent}}+\underbrace{\widehat v_F-c}_{\text{hyperscaler scarcity rent}}.\end{equation}

A leading laboratory can therefore earn both a differential return in the compute market and a downstream markup over the resulting compute-inclusive marginal cost. The precise division depends on downstream substitution, the number of laboratories, the concentration of hyperscaler capacity, and the bargaining or auction mechanism used to allocate GPUs.

\subsection{Purchaser-facing price under complete dominance}

For a simple special case, suppose one laboratory serves the market through a single leading model $L$ and receives all available capacity. Its output is

\begin{equation}y_L=\frac{K}{a_L}.\end{equation}

Let its inverse demand curve be

\begin{equation}P_L=\mathcal P_L(y_L;Q_L).\end{equation}

Assume that available capacity is below $y_L^{\mathrm{mon}}$, the laboratory's unconstrained monopoly-optimal output (the output it would choose absent the capacity constraint):

\begin{equation}\frac{K}{a_L}<y_L^{\mathrm{mon}}.\end{equation}

Under this assumption, the capacity constraint binds, and the purchaser-facing price is

\begin{equation}P_L^*=\mathcal P_L\left(\frac{K}{a_L};Q_L\right).\end{equation}

In the simple purchaser model, demand for a single active model is

\begin{equation}D_L(P_L)=N\cdot\left[1-F\left(\frac{P_L}{Q_L}\right)\right].\end{equation}

Equating demand to available output gives

\begin{equation}P_L^*=Q_L\cdot F^{-1}\left(1-\frac{K}{a_L\cdot N}\right).\end{equation}

Equivalently, the normalized purchaser-facing price per FLOP is

\begin{equation}p_L^*=q_L\cdot F^{-1}\left(1-\frac{K}{a_L\cdot N}\right).\end{equation}

The price is determined by the willingness to pay of the marginal purchaser who can be served with available capacity. It may therefore be far above physical inference cost. The dominant laboratory serves the highest-value portion of demand permitted by the capacity constraint, not all potential demand.

\subsection{Welfare: transfers, deadweight loss, and cross-laboratory misallocation}
\label{sec:welfare}

The three-component decomposition invites a welfare reading, because the components differ in whether they are transfers or genuine efficiency losses. The physical production cost $b_\ell+a_\ell\cdot c$ is a real resource cost and raises no welfare question. The hyperscaler scarcity rent $a_\ell\cdot(r_\ell-c)$ is, holding installed capacity $K$ fixed, a pure transfer from purchasers to hardware owners: with capacity binding the same quantity of output is produced whatever the level of the rent, so no deadweight loss arises from the rent itself. This transfer is nonetheless not welfare-irrelevant over a longer horizon, because it is precisely the signal that rewards capacity investment; the scarcity rent that looks like a static transfer is the dynamic incentive to build the GPUs that later relax the constraint (Section~\ref{sec:model-race}). The model-provider markup $\mu_\ell$ is the component that causes deadweight loss in the usual sense: when capacity is slack, each laboratory restricts output below the level at which price equals compute-inclusive marginal cost, and the foregone trades are a standard welfare loss (Appendix~\ref{app:slack-capacity} derives that slack-regime equilibrium explicitly).

Under binding capacity the markup's welfare consequence changes form rather than disappearing. Total output is pinned by $K$, so the markup cannot reduce aggregate quantity; what it distorts instead is the allocation of compute across laboratories. Efficient use of scarce compute equalizes marginal social value per FLOP, but oligopolistic laboratories bid for compute according to their private marginal value $\widehat v_\ell$, which nets out the cannibalization of their own inframarginal output. When markups differ across laboratories---because they face different residual elasticities---the wedge between $\widehat v_\ell$ and social marginal value differs across laboratories too, so compute is drawn toward laboratories with smaller markups relative to the surplus-maximizing allocation. The deadweight loss of downstream market power therefore does not vanish when capacity binds; it is converted from an output-reduction loss into a cross-laboratory misallocation loss, in which the wrong models are served rather than too few outputs served in total. This is the allocative counterpart, on the compute-allocation margin, of the neutrality result of Proposition~\ref{prop:capacity-neutrality}: adding laboratories does not change total output, but heterogeneous markups still move compute away from its highest-value uses.

\section{Trailing Laboratories}
\label{sec:trailing-labs}

Let $r_j$ denote trailing laboratory $j$'s actual contractual payment per FLOP, and let $r^*$ denote the current market value or opportunity cost of a FLOP. These may differ when compute was prepaid, obtained under a long-term contract, or cannot be transferred or resold.

A trailing laboratory makes a positive current operating contribution when

\begin{equation}\frac{MR_j-b_j}{a_j}\geq r_j.\end{equation}

It earns non-negative economic profit relative to the current market value of compute only when

\begin{equation}\frac{MR_j-b_j}{a_j}\geq r^*.\end{equation}

A near-frontier laboratory may satisfy this condition by charging a modest discount, serving a differentiated segment, or operating a model with greater productivity on particular tasks. A far-behind laboratory with

\begin{equation}v_j(0)<r^*\end{equation}

cannot profitably purchase even its first unit of compute at the market price. It must specialize, obtain subsidized or favorably contracted compute, use hardware with a lower opportunity cost, or leave the undifferentiated general-purpose inference market.

Subscription pricing provides an additional mechanism through which trailing laboratories may remain active without displacing the frontier provider. A purchaser may pay fixed fees to several providers in exchange for discounted future usage, and then route each task to the provider offering the greatest net value. A trailing laboratory need not become the purchaser's primary provider; it need only generate sufficient incremental value on some subset of tasks to justify its subscription fee. Appendix~\ref{app:subscriptions} formalizes this multi-homing mechanism, showing that it arises even under complete information and that it separates a provider's subscription share, usage share, and compute share, so that a trailing provider can hold a broad subscriber base while drawing only a small fraction of inference traffic.

Three forms of discounted operation should be distinguished.

\paragraph{Below average total cost.} A laboratory may fail to recover historical training costs and organizational overhead while pricing above current marginal cash cost:

\begin{equation}P_j>b_j+a_j\cdot r_j.\end{equation}

Each sale still contributes toward fixed costs.

\paragraph{Below economic opportunity cost.} A laboratory with sunk or non-transferable compute may satisfy

\begin{equation}r_j+\frac{b_j}{a_j}<p_j<r^*+\frac{b_j}{a_j}.\end{equation}

It then earns positive current cash flow but less than the value of the compute in its best alternative use. This distinction matters only when that alternative value cannot be fully realized.

\paragraph{Below marginal cash cost.} If

\begin{equation}p_j<r_j+\frac{b_j}{a_j},\end{equation}

then each additional output creates a current cash loss. Such loss-leading is rational only when continued operation generates sufficiently valuable future benefits, such as customer retention, deployment knowledge, useful data, reputation, distribution, or an increased probability of returning to the frontier. Appendix~\ref{app:dynamic-loss} formalizes this condition as a dynamic program in which operating at a current cash loss is optimal only if the discounted gain in the laboratory's continuation value exceeds the current loss; that appendix also stresses that these future benefits should not be assumed.

All of the conditions in this section are flow conditions that take the laboratory's existence---and its sunk training investment---as given. Section~\ref{sec:natural-oligopoly} asks the prior stock question: how many laboratories can expect to recover that investment at all.

\section{Dynamic Model Development and Temporary Leadership}
\label{sec:model-race}

The static model can be nested inside a dynamic race in which laboratories allocate compute and capital to building future models while serving current inference demand. Let $x_{\ell t}$ be inference FLOPs used by laboratory $\ell$ and let $g_{\ell t}$ be training or model-development FLOPs. The current aggregate constraint is

\begin{equation}\label{eq:agg-constraint}\sum_\ell x_{\ell t}+\sum_\ell g_{\ell t}\leq K_t.\end{equation}

Current inference output is

\begin{equation}y_{\ell t}=\frac{x_{\ell t}}{a_{\ell t}}.\end{equation}

Training therefore has an immediate opportunity cost: a FLOP used to build the next model cannot simultaneously serve current demand. The shadow price of current capacity applies to both uses, unless training and inference rely on non-substitutable hardware.

\subsection{Hyperscaler capacity investment}

Let $H_{ht}$ be capital expenditure by hyperscaler $h$. Because data centers and accelerators require construction and deployment time, current expenditure expands usable capacity only after a lag $T_{\mathrm{build}}$:

\begin{equation}K_{h,t+1}=\left(1-\delta_K\right)\cdot K_{ht}+\Phi_h\left(H_{h,t-T_{\mathrm{build}}}\right).\end{equation}

Aggregate capacity is $K_t=\sum_h K_{ht}$. If an additional unit of installed capacity earns future scarcity rent $\lambda_{h,t+s}$, the hyperscaler's investment condition is

\begin{equation}C_h'(H_{ht})=\mathbb E_t\left[\sum_{s=T_{\mathrm{build}}}^{\infty}\beta^s\cdot\lambda_{h,t+s}\cdot\frac{\partial K_{h,t+s}}{\partial H_{ht}}\right].\end{equation}

Expected future model improvements can therefore raise current GPU and infrastructure investment by increasing expected future demand and the shadow value of capacity. Capacity expansion reduces scarcity conditional on demand, but rising model value and training demand can outpace the construction of new capacity. The rents $\lambda_{h,t+s}$ in this condition are a binding-regime object: they fall continuously to zero as capacity approaches the slack threshold of Appendix~\ref{app:slack-capacity}, which is what keeps forward-looking investment from planning capacity past that boundary; the appendix works out this implication.

\subsection{Lumpy model releases}

Let $R_{\ell t}\in\{0,1\}$ indicate whether laboratory $\ell$ completes and releases a new model. The release probability depends on training compute and other research investment $I_{\ell t}$:

\begin{equation}\Pr\left(R_{\ell t}=1\mid g_{\ell t},I_{\ell t}\right)=h_\ell(g_{\ell t},I_{\ell t}),\qquad \frac{\partial h_\ell}{\partial g_{\ell t}}>0,\qquad \frac{\partial h_\ell}{\partial I_{\ell t}}>0.\end{equation}

Conditional on a release, the laboratory draws a candidate productivity level

\begin{equation}\log q_{\ell,t+1}^{\mathrm{new}}=A_{t+1}+\alpha_\ell+\nu_{\ell,t+1}.\end{equation}

The deployed productivity becomes

\begin{equation}q_{\ell,t+1}=\begin{cases}\max\left\{q_{\ell t},q_{\ell,t+1}^{\mathrm{new}}\right\},&R_{\ell t}=1,\\q_{\ell t},&R_{\ell t}=0.\end{cases}\end{equation}

The static market equilibrium at date $t$ maps the deployed productivity vector and available inference capacity into target compute allocations:

\begin{equation}\widehat{\mathbf x}_t=\mathbf x^*(\mathbf q_t,K_t^I),\qquad K_t^I\equiv K_t-\sum_\ell g_{\ell t}.\end{equation}

A successful release shifts the laboratory's private marginal-profit schedule upward and can cause a discrete increase in its target compute and market shares. To represent switching, evaluation, and integration delays, actual compute share may adjust gradually:

\begin{equation}s_{\ell,t+1}=\left(1-\sigma\right)\cdot s_{\ell t}+\sigma\cdot\widehat s_{\ell t},\qquad \widehat s_{\ell t}\equiv\frac{\widehat x_{\ell t}}{K_t^I},\qquad \sigma\in(0,1].\end{equation}

\subsection{Temporary leadership rents}

Let $\pi_L^{\mathrm{lead}}$ be the leader's flow profit and $\pi_L^{\mathrm{follow}}$ its flow profit after losing the lead. Define

\begin{equation}\Delta\pi_L\equiv\pi_L^{\mathrm{lead}}-\pi_L^{\mathrm{follow}}.\end{equation}

Let $\Lambda_L$ be the hazard that a competitor releases a sufficiently strong substitute. With discount rate $\rho$ (distinguished here from the wholesale compute price $r$, $r^*$ of Sections~\ref{sec:concentration}--\ref{sec:oligopoly}, which plays no role in this formula), the approximate value of temporary leadership is

\begin{equation}\boxed{V_L^{\mathrm{lead}}\simeq\frac{\Delta\pi_L}{\rho+\Lambda_L}.}\end{equation}

The formula is the standard exponential-hazard perpetuity. Over a short interval $dt$ the leader collects the flow premium $\Delta\pi_L\,dt$ and keeps the lead with probability $1-\Lambda_L\,dt$, so $V_L^{\mathrm{lead}}=\Delta\pi_L\,dt+(1-\rho\,dt)(1-\Lambda_L\,dt)\,V_L^{\mathrm{lead}}$; discarding the second-order term and letting $dt\to0$ gives the boxed expression. Three assumptions are hidden in it: (i) the flow premium $\Delta\pi_L$ is constant while the lead lasts; (ii) the catch-up hazard $\Lambda_L$ is constant, so displacement is a Poisson event; and (iii) losing the lead is absorbing. The last is the strongest. Strictly, $V_L^{\mathrm{lead}}$ is the difference between the lead-state and follow-state value functions, and the follow state itself carries an option to regain the lead; the formula sets the value of that option to zero and so overstates the cost of being displaced. A fourth restriction is regime-dependence: the flow premium $\Delta\pi_L$ bundles a differential compute rent and a downstream markup, and in the slack-capacity regime of Appendix~\ref{app:slack-capacity} the compute-rent component is zero, so the formula still applies but the premium is built from the markup channel alone.

The approximation is also partial-equilibrium in $\Lambda_L$. The hazard is written here as an exogenous constant, but the release process above makes hazards endogenous through $h_\ell(g_{\ell t},I_{\ell t})$, and the investment condition below implies that followers invest more when $V_L^{\mathrm{lead}}$ is large. A higher leadership value therefore raises follower investment and hence $\Lambda_L$, which feeds back to lower $V_L^{\mathrm{lead}}$. So $\Lambda_L$ is an equilibrium object rather than a free parameter: the two should not be varied as independent dials, and the formula is best read as a fixed point evaluated at the equilibrium hazard. We do not solve this fixed point here.

A follower invests because a successful release may move it from thin-margin or negligible demand to temporary leadership. Because the follower would replace another laboratory's rents while the leader would largely cannibalize its own, the follower's private gain from a given release exceeds the leader's, which is the replacement effect of \citet{arrow1962}. The leader nonetheless invests both to advance the frontier and to raise the threshold that a follower must cross; this second, defensive motive is the preemption incentive identified by \citet{gilbertNewbery1982}. The investment problem is a race of the kind studied by \citet{loury1979}, \citet{leeWilde1980}, and \citet{reinganum1983}, in which each laboratory spends to raise its own release hazard while facing the hazard generated by rivals.

Two first-order conditions make the cost side explicit, because the release hazard $h_\ell(g_{\ell t},I_{\ell t})$ depends on both training compute $g_{\ell t}$ and other research investment $I_{\ell t}$. For research investment,

\begin{equation}C_\ell'(I_{\ell t})=\frac{\partial h_\ell}{\partial I_{\ell t}}\cdot\mathbb E_t\left[\left(V_{\ell,t+1}^{\mathrm{new}}-V_{\ell,t+1}^{\mathrm{old}}\right)_+\mid R_{\ell t}=1\right].\end{equation}

For training compute the marginal cost is not a direct outlay but the forgone inference margin emphasized at the start of this section: a FLOP allocated to $g_{\ell t}$ cannot serve current demand, so it costs its cash cost $c$ to run plus the shadow rent $\lambda_t$ on the scarce capacity it occupies, where $\lambda_t$ is the multiplier on the aggregate constraint~\eqref{eq:agg-constraint}. The condition on training compute is therefore

\begin{equation}c+\lambda_t=\frac{\partial h_\ell}{\partial g_{\ell t}}\cdot\mathbb E_t\left[\left(V_{\ell,t+1}^{\mathrm{new}}-V_{\ell,t+1}^{\mathrm{old}}\right)_+\mid R_{\ell t}=1\right].\end{equation}

Writing the training-compute condition against the full opportunity cost $c+\lambda_t$, rather than treating $g_{\ell t}$ as free, keeps the race consistent with the section's opening observation that training and inference compete for the same scarce capacity. The same opportunity-cost price per training FLOP is what converts cumulative training compute into a monetary entry cost in Section~\ref{sec:natural-oligopoly}. The private return in both conditions includes both new social value and business stealing from competitors. Model-race investment can therefore be excessive because laboratories duplicate costly training efforts to capture temporary rents, or insufficient because they do not capture all consumer and downstream productivity benefits.

\section{Technological Regimes, Consumer Prices, and Rent Shares}
\label{sec:technology-regimes}

Let $G_{\ell t}$ denote cumulative effective training investment by laboratory $\ell$, and let productivity per inference FLOP be $q_{\ell t}=f_\ell(G_{\ell t})$. A notational note before proceeding: this $G_{\ell t}$ (cumulative training compute, indexed by time) is unrelated to the gross operating value function $G_\ell(x_\ell)$ of Section~\ref{sec:concentration} (a function of a compute allocation, not a stock); the collision runs through the rest of this section and Section~\ref{sec:natural-oligopoly}, where $\Gamma_\ell$ is introduced specifically to avoid compounding it further (Section~\ref{subsec:entry-cost}). Context disambiguates the two uses of $G$ throughout---the operating-value function always takes a compute allocation as its argument, the training-investment stock never does---but both are flagged here in one place.

The long-run market dynamics depend on the asymptotic behavior of the training-to-quality map $q_\ell(G)$, and the cases are organized by two questions that are easy to conflate. The first is about \emph{levels}: does $q_\ell(G)$ approach a ceiling as cumulative training investment grows, or improve without bound? The second is about \emph{relative position}: does the gap $q_L(G)/q_F(G)$ converge to one, or stay bounded above it? The two questions are independent, and it is the second, not the first, that governs whether differential productivity rents survive, because the machinery of Sections~\ref{sec:persistent-advantage}--\ref{sec:oligopoly} prices a laboratory's advantage through its quality relative to rivals'. (More precisely, convergence of relative position guarantees that differential rents vanish \emph{as a share} of the value the market generates; whether they also vanish in level turns on the absolute gap, and Section~\ref{subsec:additive-frontier} works out the one cell of the table in which the two versions of the statement come apart.) Crossing the two questions gives four asymptotic regimes, shown in Table~\ref{tab:regime-grid} and taken up in turn below.

\begin{table}[htbp]
\centering
\caption{The four technological regimes, organized by whether quality levels stay bounded and whether the relative gap closes.}
\label{tab:regime-grid}
\begin{tabular}{lll}
\hline
 & Gap $q_L/q_F\rightarrow1$ & Gap persists \\
\hline
Levels bounded & Common asymptote (Sec.~\ref{subsec:common-asymptote}) & Firm-specific asymptotes (Sec.~\ref{subsec:firm-specific-asymptotes}) \\
Levels unbounded & Additive advantage (Sec.~\ref{subsec:additive-frontier}) & Multiplicative advantage (Sec.~\ref{subsec:continuing-improvement}) \\
\hline
\end{tabular}
\end{table}

One clarification about the first question before proceeding: $q$ is economically relevant quality per FLOP---denominated in purchaser value, not in benchmark scores---so demand-side saturation is not a separate technological future: an ``improvement'' that no longer raises what purchasers value is not an increase in $q$ at all, and saturating demand for marginal quality therefore simply \emph{is} a ceiling in $q$, landing in the bounded row of Table~\ref{tab:regime-grid} in the common- or firm-specific-asymptote case according to whether laboratories reach the point of purchaser indifference together or at different levels.

Under oligopoly, the purchaser-facing price is

\begin{equation}P_{\ell t}=b_{\ell t}+a_{\ell t}\cdot r_{\ell t}+\mu_{\ell t}.\end{equation}

The first two terms are compute-inclusive marginal cost; $\mu_{\ell t}$ is the downstream oligopoly markup. When upstream compute is allocated competitively, $r_t$ depends on the best alternative laboratory's private marginal value. When laboratories exercise oligopsony or bargain bilaterally, the division is modified as described in Appendix~\ref{app:bargaining}.

\subsection{A common global asymptote}
\label{subsec:common-asymptote}

Suppose all laboratories approach the same maximum productivity $\bar q$:

\begin{equation}q_\ell(G)=\bar q-d_\ell\cdot(\psi_\ell\cdot G)^{-\beta},\qquad d_\ell>0,\qquad \psi_\ell>0,\qquad \beta>0.\end{equation}

Here $\psi_\ell>0$ is a laboratory-specific research-efficiency (rate-of-approach) parameter, distinct from the productivity-level multiplier $\eta_\ell$ of Section~\ref{sec:persistent-advantage}: it governs how fast laboratory $\ell$ closes the gap to the common ceiling, but has no effect on the limit. The remaining distance to the asymptote is

\begin{equation}\bar q-q_\ell(G)=d_\ell\cdot(\psi_\ell\cdot G)^{-\beta}.\end{equation}

Reducing this gap by a constant proportional factor requires a multiplicative increase in training compute. After $n$ equal proportional reductions, required compute grows exponentially in $n$.

Despite differences in research efficiency, all laboratories converge to the same productivity:

\begin{equation}\lim_{G\rightarrow\infty}q_\ell(G)=\bar q,\qquad \lim_{G\rightarrow\infty}\frac{q_L(G)}{q_F(G)}=1.\end{equation}

If downstream products become close substitutes and laboratories compete in prices, convergence raises own-price elasticities and tends to compress model-provider markups:

\begin{equation}\epsilon_{\ell\ell}\uparrow\;\Longrightarrow\;\mu_\ell\downarrow.\end{equation}

This chain is the quality-convergence direction of Proposition~\ref{prop:breakthrough} (stated below in Section~\ref{subsec:continuing-improvement}): in the two-model demand system it is exactly the statement that a shrinking quality gap raises the leader's own-price elasticity and compresses its markup. However, technological convergence does not by itself imply competitive pricing. With a finite Cournot oligopoly, collusion, switching costs, or ecosystem lock-in \citep{klemperer1995}, positive model-provider markups may persist even when model qualities converge.

Training investment eventually weakens because the marginal productivity return approaches zero. If hyperscaler capacity continues to expand, the wholesale scarcity component $r_t-c_t$ also tends to fall. The strongest commoditization outcome therefore requires both technological convergence and sufficient downstream and upstream competition:

\begin{equation}q_L-q_F\rightarrow0,\qquad \mu_\ell\rightarrow0,\qquad r_t-c_t\rightarrow0.\end{equation}

Under these conditions, quality-adjusted consumer prices approach physical inference cost; this commoditization endpoint is the slack-capacity regime of Appendix~\ref{app:slack-capacity} combined with vanishing differentiation. If downstream oligopoly persists while compute becomes abundant, hyperscaler rents disappear but laboratories retain markups. If compute remains scarce, hyperscalers may capture a large share of a shrinking model-specific surplus even as absolute scarcity rents remain positive.

\subsection{Firm-specific asymptotes}
\label{subsec:firm-specific-asymptotes}

Suppose instead that laboratories approach different structural productivity limits:

\begin{equation}q_\ell(G)=\eta_\ell\cdot\bar q-d_\ell\cdot G^{-\beta}.\end{equation}

Then

\begin{equation}\lim_{G\rightarrow\infty}\frac{q_L(G)}{q_F(G)}=\frac{\eta_L}{\eta_F}>1.\end{equation}

Training returns still taper, but product differentiation does not disappear. The leader retains a persistent shift in both its demand and private marginal-profit schedules. Complete dominance is possible if

\begin{equation}\widehat v_L(K;\eta_L\cdot\bar q)\geq\widehat v_F(0;\eta_F\cdot\bar q),\end{equation}

while weaker differences produce a vertically differentiated oligopoly.

Consumer prices retain a durable frontier premium. Even if compute becomes abundant so that $r_t\rightarrow c_t$, the leader may charge a positive quality and market-power markup:

\begin{equation}P_L-\left(b_L+a_L\cdot c\right)=\mu_L>0.\end{equation}

If compute remains scarce, purchaser prices contain both the hyperscaler scarcity component and the downstream model-provider markup. The leading laboratory retains a larger share of producer surplus than under the common-asymptote regime because both differential productivity and downstream differentiation remain positive.

\subsection{Continuing improvement with multiplicative advantage}
\label{subsec:continuing-improvement}

Suppose there is no nearby global ceiling. A simple specification is

\begin{equation}q_{\ell t}=\eta_\ell\cdot\bar q_t\cdot\xi_{\ell t},\qquad \bar q_{t+1}=\left(1+g_t\right)\cdot\bar q_t,\end{equation}

or, as a function of cumulative effective training compute,

\begin{equation}q_\ell(G)=\eta_\ell\cdot G^\beta.\end{equation}

The multiplicative form is itself a choice among the unbounded cases, not forced by unboundedness: it makes the leader's proportional advantage $\eta_L/\eta_F$ independent of the frontier level, so the gap that drives the release-cycle dynamics below never trends away. The additive alternative, in which the frontier still diverges but relative positions converge, is the fourth regime and is taken up in Section~\ref{subsec:additive-frontier}.

A major release increases the leader's product quality and reduces substitutability. Two distinct downstream consequences follow, and they are established in two distinct models of price and quantity competition, because the paper's maintained binding-capacity regime and the regime in which unconstrained price competition is well posed are not the same setting (Section~\ref{sec:limitations} returns to this scope restriction). Immediately after a breakthrough, in the slack-capacity regime where compute is obtainable at the going compute-inclusive cost,

\begin{equation}q_L\uparrow\;\Longrightarrow\;\epsilon_{LL}\downarrow\;\Longrightarrow\;\mu_L\uparrow,\end{equation}

and, separately, in the capacity-constrained regime where the leader holds the entire binding compute endowment,

\begin{equation}q_L\uparrow\;\Longrightarrow\;\widehat v_L-\widehat v_F\uparrow.\end{equation}

The leader may harvest the breakthrough through higher prices and restricted output, or maintain a lower price to acquire subscribers, integrations, and switching costs before competitors catch up \citep{klemperer1987}. The robust prediction is therefore not that nominal prices always rise, but that the leader's feasible price-quality frontier, potential markup, and differential compute rent all increase.

When a competitor releases a near-frontier model, substitution increases and both chains reverse:

\begin{equation}q_F\uparrow\;\Longrightarrow\;\epsilon_{LL}\uparrow\;\Longrightarrow\;\mu_L\downarrow,\qquad q_F\uparrow\;\Longrightarrow\;\widehat v_L-\widehat v_F\downarrow,\end{equation}

and the hyperscaler's outside option improves. These arrow chains are not free-floating assertions about an unspecified demand system: in a general demand system quality convergence need not raise equilibrium own-price elasticities, since prices adjust as well, and markups depend on curvature rather than elasticity alone; nor need the differential value of compute move monotonically with the gap in an arbitrary capacity-allocation mechanism. Both chains do hold, however, in the two-model demand system of Lemma~\ref{lem:two-model}---the first under Bertrand competition in the linear workhorse of \citet{singhVives1984}, the second under capacity-constrained Cournot competition in the same demand system, which is the reduced form for capacity-then-price competition due to \citet{krepsScheinkman1983} cited in Section~\ref{sec:literature-review}---which lets us state the headline dynamic prediction as a proposition rather than a claim.

\begin{proposition}[Breakthroughs, catch-up, and the division of surplus]\label{prop:breakthrough}
Let $\Delta\equiv Q_L-Q_F>0$ be the quality gap in the two-model demand system of Lemma~\ref{lem:two-model}. Two comparative statics are established, each in the model regime in which it is actually well posed; Appendix~\ref{app:two-model} gives both derivations.
\begin{enumerate}
\item \textbf{Markup and elasticity (Bertrand, slack capacity).} Under Bertrand price competition with a common compute-inclusive marginal cost $\kappa\geq0$ obtainable at the going wholesale price---i.e., in the slack-capacity regime of Appendix~\ref{app:slack-capacity}, where nothing rations either laboratory's output below its unconstrained price-competition choice---and $\kappa<2Q_L$ (equivalently, an interior equilibrium in which the leader's markup is positive), the leader's equilibrium Lerner markup $\mu_L=\Delta(2Q_L-\kappa)/(3Q_L+\Delta)$ is strictly increasing in $\Delta$ holding $Q_L$ and $\kappa$ fixed, with $\mu_L\to0$ as $\Delta\to0$. Its equilibrium own-price elasticity $\epsilon_{LL}=P_L/\mu_L$ is strictly decreasing in $\Delta$ whenever $\kappa>0$; at $\kappa=0$, $\epsilon_{LL}\equiv1$ for every $\Delta$, a degenerate special case rather than the general result.
\item \textbf{Differential compute rent (capacity-constrained Cournot, binding capacity).} Under aggregate capacity-constrained Cournot competition---each laboratory's compute endowment $x_\ell$ enters explicitly, and price competition subject to that endowment reduces to this quantity game by the Kreps--Scheinkman logic cited above---evaluated at the winner-take-all corner $x_L=K$, $x_F=0$ with $a_L=a_F=1$ and common compute-inclusive cost $\kappa_L=\kappa_F$, the differential private value of compute is $\widehat v_L-\widehat v_F=-K\,Q_L+\Delta\,(1-K)$, strictly increasing in $\Delta$ holding $Q_L$ and $\kappa$ fixed, whenever $K<1$---the aggregate scarcity condition maintained throughout the paper (Section~\ref{sec:fixed-capacity}).
\end{enumerate}
Consequently a breakthrough that widens the gap lowers $\epsilon_{LL}$ and raises $\mu_L$ (part~1) and $\widehat v_L-\widehat v_F$ (part~2); a follower release that narrows the gap reverses both. Part~2 holds the compute-inclusive cost $\kappa$ fixed as $\Delta$ varies; it is a ceteris-paribus, partial-equilibrium statement about the immediate effect of a quality change, not a claim that $\kappa$ (through $r^*$) stays fixed once the market re-equilibrates around a follower's release---see the discussion following the proof in Appendix~\ref{app:two-model} for why the partial-equilibrium version remains informative despite this.
\end{proposition}

Laboratories therefore tend to capture a larger share of producer surplus immediately after breakthroughs, while hyperscalers tend to capture a larger share as competitors catch up. Absolute hyperscaler profits may rise in both phases if model improvements expand demand for inference and training compute. The proposition concerns the division of producer surplus, not necessarily the level of profit.

\subsection{Continuing improvement with additive advantage}
\label{subsec:additive-frontier}

The remaining cell of Table~\ref{tab:regime-grid} is unbounded improvement with a converging relative position, and the paper already contains the specification that produces it: the additive alternative that Section~\ref{sec:persistent-advantage} introduced and deliberately set aside---set aside precisely because a fixed additive increment cannot represent \emph{persistent} advantage against an advancing frontier. That earlier observation was not a dead end; it is the definition of this regime. Suppose

\begin{equation}q_\ell(G)=\bar q(G)+\alpha_\ell,\qquad \bar q(G)=G^{\beta}\rightarrow\infty,\end{equation}

or in time-series form $q_{\ell t}=\bar q_t+\alpha_\ell$ with $\bar q_t\rightarrow\infty$, where $\alpha_\ell$ is the additive laboratory increment of Section~\ref{sec:persistent-advantage}'s additive specification. Levels diverge, the absolute gap is frozen at $q_L-q_F=\alpha_L-\alpha_F$, and the relative gap closes:

\begin{equation}\lim_{G\rightarrow\infty}q_\ell(G)=\infty,\qquad \lim_{G\rightarrow\infty}\frac{q_L(G)}{q_F(G)}=1.\end{equation}

The second limit is exactly the one computed in Section~\ref{sec:persistent-advantage}, $\lim_{A_t\rightarrow\infty}(A_t+\alpha_L)/(A_t+\alpha_F)=1$; this regime is where that calculation pays off. There, the vanishing ratio was the reason to reject the additive form as a representation of persistent advantage. Here it is not a modeling defect but a candidate description of the world: a frontier that keeps advancing while every laboratory's head start becomes proportionally negligible.

There is no commoditization by ceiling in this regime. The marginal product of training compute behaves exactly as in the multiplicative case---returns taper at the rate set by $\beta$, not toward zero at an asymptote---so training investment is not extinguished by exhausted technology, quality keeps rising without bound, and quality-adjusted consumer prices can keep falling indefinitely through the productivity channel rather than converging to the positive floor that physical inference cost and a finite ceiling $\bar q$ jointly imply in the bounded regimes. What erodes is instead the leader's position relative to the frontier: commoditization by convergence, without commoditization by ceiling.

The two-model demand system makes the erosion precise, and also shows that it is subtler than the common-asymptote case. The rent-bearing objects of Proposition~\ref{prop:breakthrough} depend on the \emph{absolute} quality gap $\Delta=Q_L-Q_F$, and this regime pins that gap: with $a_L=a_F=1$, $\Delta=\alpha_L-\alpha_F$ at every frontier level. Substituting a diverging $Q_L$ at fixed $\Delta$ into the proposition's expressions gives the regime's signature, derived within the linear example. The Bertrand markup of part~1 tends to a positive constant, $\mu_L=\Delta(2Q_L-\kappa)/(3Q_L+\Delta)\rightarrow\tfrac{2}{3}\Delta$, so the leader's price stays bounded while $Q_L\rightarrow\infty$ and the quality-adjusted price $P_L/Q_L\rightarrow0$. The differential rent therefore does not vanish in \emph{level}: it converges to an absolute margin pinned by the additive increment alone, independent of how far the frontier has advanced. It vanishes in \emph{share}: relative to purchaser value, consumer surplus, and the revenue the market generates---all of which grow with the frontier---the frozen margin becomes negligible. This is the one cell of Table~\ref{tab:regime-grid} in which the share and level versions of ``differential rents vanish'' come apart, which is why the section's framing insists that convergence of relative position governs rent \emph{shares}: under the common asymptote the absolute gap closes too, so rents vanish in level and in share at once, whereas here only the share vanishes. By the same substitution, the equilibrium elasticity $\epsilon_{LL}=P_L/\mu_L$ tends to the constant $1+3\kappa/(2\Delta)$ rather than diverging, so ratio convergence alone does not deliver the homogeneous-goods limit; that requires the absolute gap to close as well.

The capacity-allocation side is starker. At the winner-take-all corner of Proposition~\ref{prop:breakthrough}(2), the differential compute rent is $\widehat v_L-\widehat v_F=-K\,Q_L+\Delta\,(1-K)$; at fixed $\Delta$ this turns negative once $Q_L>\Delta\,(1-K)/K$, so continued frontier growth eventually breaks the winner-take-all condition of Section~\ref{sec:concentration} by itself---the leader's marginal FLOP, serving ever-lower-value purchasers with ever-higher quality, comes to be worth less than the follower's first FLOP serving the top of the market at a quality only $\Delta$ lower. Under multiplicative advantage the same expression is scale-invariant: there $\Delta=(1-\eta_F/\eta_L)\cdot Q_L$, so $\widehat v_L-\widehat v_F=Q_L\cdot\left[(1-\eta_F/\eta_L)(1-K)-K\right]$ keeps a fixed sign at every frontier level. Whether concentration is sustainable is thus level-independent under multiplicative advantage and self-extinguishing under additive advantage. Both statements are corner evaluations in the linear example, offered as illustrations of the mechanism rather than general theorems.

Two carry-overs from the other regimes complete the picture. As under the common asymptote, convergence removes the vertical-differentiation source of market power but not the structural ones: a finite oligopoly, collusion, switching costs, or ecosystem lock-in \citep{klemperer1995} can sustain model-provider markups even after the differentiation-based component has shrunk to the frozen $\Delta$-margin, so this regime no more guarantees competitive pricing than the common asymptote does. And as under multiplicative continuing improvement, temporary model-specific shocks and lumpy releases still open transient gaps along the path, so the breakthrough-and-catch-up dynamics of Proposition~\ref{prop:breakthrough} apply unchanged period by period; what the regime removes is the trend---leadership episodes recur, but no laboratory's proportional position survives them.

Which of the two unbounded regimes describes reality is exactly the additive-versus-multiplicative question that Section~\ref{sec:persistent-advantage} posed when choosing how to represent persistent advantage, now revealed as a question about long-run market structure rather than about notation: a multiplicative advantage makes continuing improvement a sequence of durable proportional leads, while an additive advantage makes it a convergence machine. The observable difference is whether the leader's proportional lead is stable or shrinking as the frontier advances. The paper takes no stand on which obtains: the multiplicative form was adopted in Section~\ref{sec:persistent-advantage} because it is the only one of the two that can represent persistent advantage at all, not as a claim that the additive world is less likely.

\subsection{Comparison of the regimes}
\label{subsec:regime-comparison}

The key distinction is not simply whether performance eventually stops improving. It is whether the leader--follower gap and downstream market power disappear---the second axis of Table~\ref{tab:regime-grid}, which cuts across the first.

Under a common asymptote, differential productivity rents vanish---in level and in share, since the absolute and relative gaps close together---but oligopoly markups can persist if competition remains weak. Hyperscaler rents vanish only if capacity becomes abundant or upstream competition drives the compute price to physical cost.

Under firm-specific asymptotes, the leader retains both a durable productivity advantage and potentially a durable downstream markup. Hyperscalers capture the scarcity value common to viable laboratories, but the leader retains a positive differential rent.

Under continuing improvement with multiplicative advantage, market shares and rent shares fluctuate around model releases. The leader captures a larger portion of surplus after a breakthrough; competitors and hyperscalers gain bargaining power during catch-up. Consumer prices can rise, remain sticky, or fall in nominal terms depending on laboratory strategy, but quality-adjusted prices fall only when productivity and capacity growth outpace demand and oligopoly markups.

Under continuing improvement with additive advantage, the release-cycle fluctuations are the same, but the trend belongs to the convergent column: differential rents freeze in level and vanish as a share of a market whose total value keeps growing, winner-take-all becomes unsustainable at high frontier levels, and quality-adjusted prices can fall indefinitely because there is no ceiling to arrive at. The regime shares its convergence logic with the common asymptote and its unbounded growth with the multiplicative case, and matches neither: rents do not vanish in level as under the ceiling, and they do not persist in share as under multiplicative advantage.

The four regimes are asymptotic classifications, not forecasts of a single uninterrupted path: to the extent the world moves between them over time---a stretch of continuing improvement that later runs into a ceiling, say---each phase is analyzed by applying the regime that describes it, in sequence, and no fifth case is needed.

\section{Fixed Training Costs and the Number of Laboratories}
\label{sec:natural-oligopoly}

Everything so far takes the set of laboratories as given and asks how they divide scarce compute, how they price, and how rents split between them and the hyperscalers. This section asks the question upstream of all of that: how many laboratories can exist in the first place? Training a frontier model is a large, sunk, essentially fixed cost, while serving additional inference once the model exists is comparatively cheap at the margin. That combination is the classic setup for a natural monopoly or natural oligopoly: the number of firms is bounded not by who wins the contest over scarce compute, but by how many can expect to recover the fixed cost of showing up equipped to compete. The winner-take-all condition of Section~\ref{sec:concentration}, and its segment-level version in Proposition~\ref{prop:segment-wta}, operate on the demand side among laboratories that already hold frontier models; the free-entry condition developed here determines how many laboratories reach that stage at all.

\subsection{The fixed cost of a frontier model}
\label{subsec:entry-cost}

The model already contains the two ingredients needed to price entry, so we do not introduce a separate fixed-cost primitive. First, Section~\ref{sec:technology-regimes} relates quality per inference FLOP to cumulative effective training compute through a regime-specific law $q_\ell(G)$. Inverting it gives the training compute required to field a model of quality $q$:

\begin{equation}\Gamma_\ell(q)\equiv q_\ell^{-1}(q).\end{equation}

(We write $\Gamma_\ell$ rather than reusing $G_\ell$, which already denotes gross operating value in Section~\ref{sec:concentration}.) Second, Section~\ref{sec:model-race} establishes the price of a training FLOP: its full opportunity cost $c+\lambda_t$, the cash cost of running it plus the shadow rent on the scarce capacity it occupies. The monetary fixed cost of entering at quality $q$ is therefore

\begin{equation}\label{eq:entry-cost}\mathcal F_\ell(q)=\int_0^{\Gamma_\ell(q)}\left(c_t+\lambda_t\right)dG\;\simeq\;\left(c+\lambda\right)\cdot\Gamma_\ell(q),\end{equation}

where the approximation treats the opportunity cost of compute as roughly constant over the training program.

The four technological regimes of Section~\ref{sec:technology-regimes} give the inverse laws directly, and they say quite different things about entry:

\begin{itemize}
\item \textbf{Common asymptote.} $q_\ell(G)=\bar q-d_\ell\cdot(\psi_\ell\cdot G)^{-\beta}$ inverts to $\Gamma_\ell(q)=\psi_\ell^{-1}\cdot\left(d_\ell/(\bar q-q)\right)^{1/\beta}$, which diverges as $q\to\bar q$. Entry near the ceiling is arbitrarily expensive, so the entry cost rises fastest exactly where products are converging and markups are compressing: late entrants pay more to obtain less differentiation.
\item \textbf{Firm-specific asymptotes.} $q_\ell(G)=\eta_\ell\cdot\bar q-d_\ell\cdot G^{-\beta}$ inverts to $\Gamma_\ell(q)=\left(d_\ell/(\eta_\ell\cdot\bar q-q)\right)^{1/\beta}$ for $q<\eta_\ell\cdot\bar q$, and no finite training budget reaches $q\geq\eta_\ell\cdot\bar q$. Which laboratories can enter at a commercially required quality is then determined by ceilings before budgets: a laboratory whose asymptote lies below that quality is excluded at any price.
\item \textbf{Continuing improvement, multiplicative advantage.} $q_\ell(G)=\eta_\ell\cdot G^{\beta}$ inverts to $\Gamma_\ell(q)=\left(q/\eta_\ell\right)^{1/\beta}$, finite for every target. The persistent multiplier $\eta_\ell$ of Section~\ref{sec:persistent-advantage} now appears as an entry-cost advantage, $\mathcal F_\ell(q)\propto\eta_\ell^{-1/\beta}$, so free entry selects laboratories in descending order of research efficiency; and when $\beta<1$, required compute is convex in target quality, so each further increment to the entry bar costs proportionally more.
\item \textbf{Continuing improvement, additive advantage.} $q_\ell(G)=G^{\beta}+\alpha_\ell$ inverts to $\Gamma_\ell(q)=\left(q-\alpha_\ell\right)^{1/\beta}$, finite for every target above $\alpha_\ell$. The additive increment appears as an entry-cost discount, but a proportionally vanishing one: $\Gamma_\ell(q)/\Gamma_j(q)=\left((q-\alpha_\ell)/(q-\alpha_j)\right)^{1/\beta}\rightarrow1$ as $q\rightarrow\infty$, so laboratory identity matters less for the cost of entry the higher the entry bar rises---the entry-cost mirror of the regime's converging relative quality, and the opposite of the multiplicative case, where the discount $\eta_\ell^{-1/\beta}$ is the same proportional factor at every target quality.
\end{itemize}

One honest caveat before using $\mathcal F_\ell$. Equation~\eqref{eq:entry-cost} prices the compute component of entry, because that is what the model contains. Real fixed costs also include research talent, data acquisition, failed and duplicated training runs, and organizational overhead, none of which the model represents, so $\mathcal F_\ell$ should be read as a lower bound, and the model makes no claim to quantify the true figure. The results below need only that the fixed cost is large, sunk, and paid before any revenue arrives.

\subsection{Free entry and the equilibrium number of laboratories}
\label{subsec:free-entry}

A laboratory pays $\mathcal F$ only if the present value of the profits it expects afterward covers the outlay. The model already supplies the per-period profit of an active laboratory: in the symmetric binding-capacity benchmark of Proposition~\ref{prop:capacity-neutrality}, each of $n$ laboratories earns markup $B\cdot K/(a\cdot n)$ on output $K/(a\cdot n)$, so flow profit is

\begin{equation}\pi(n)=\frac{B\cdot K^{2}}{a^{2}\cdot n^{2}}\end{equation}

per period (Appendix~\ref{app:cournot}). Two features of this expression drive everything below. It is strictly decreasing in $n$; and it falls like $1/n^{2}$ rather than $1/n$, because each entrant both splits the laboratories' profit pool and shrinks it---the aggregate pool $B\cdot K^{2}/(a^{2}\cdot n)$ itself declines in $n$ as competition transfers fixed-capacity rent upstream to hyperscalers (Proposition~\ref{prop:capacity-neutrality}). Entry into the laboratory layer is therefore doubly deterred while capacity binds.

Discount flow profits at rate $\rho$; if a laboratory's rents also face a displacement hazard $\Lambda$, discount at the effective rate $\rho+\Lambda$, by the same hazard-perpetuity logic as the temporary-leadership value of Section~\ref{sec:model-race}. Free entry then pins down the number of laboratories: this is the endogenous-sunk-cost logic of \citet{sutton1991}, applied with training compute as the sunk outlay. The flow profit $\pi(n)$ used below is the competitive-upstream Cournot profit of Appendix~\ref{app:cournot}, so this proposition inherits the same price-taking-in-$r^*$ conduct assumption as Proposition~\ref{prop:capacity-neutrality}, with the same qualification: under Nash bargaining with laboratory weight $\zeta>0$ (Appendix~\ref{app:bargaining}), laboratories retain part of the scarcity rent that price-taking here transfers upstream, $\pi(n)$ is correspondingly larger, and the free-entry count $n^*$ below is a lower bound on the bargained-conduct count rather than an estimate of it.

\begin{proposition}[Free entry under binding capacity]\label{prop:free-entry}
Let entrants at quality parity pay a common fixed cost $\mathcal F>0$ and, once active among $n$ laboratories, earn the flow profit $\pi(n)=B\cdot K^{2}/(a^{2}\cdot n^{2})$ of Appendix~\ref{app:cournot}, discounted at effective rate $\rho+\Lambda>0$, and suppose capacity binds at the resulting number of entrants (the threshold condition of Appendix~\ref{app:cournot}), except where part~(4) explicitly considers the slack case. Then:
\begin{enumerate}
\item there is a unique free-entry number of laboratories---the largest $n$ such that $\pi(n)/(\rho+\Lambda)\geq\mathcal F$---given by
\begin{equation}n^{*}=\left\lfloor\frac{K}{a}\cdot\sqrt{\frac{B}{(\rho+\Lambda)\cdot\mathcal F}}\right\rfloor\end{equation}
whenever $\pi(1)/(\rho+\Lambda)\geq\mathcal F$; otherwise no laboratory enters;
\item the market is a natural monopoly---a single laboratory clears the bar while a second entrant would not---precisely when
\begin{equation}\frac{B\cdot K^{2}}{4\cdot a^{2}}<(\rho+\Lambda)\cdot\mathcal F\leq\frac{B\cdot K^{2}}{a^{2}},\end{equation}
a fourfold window of fixed costs, since $\pi(1)/\pi(2)=4$;
\item $n^{*}$ is weakly decreasing in $\mathcal F$, $\rho$, and $\Lambda$, and weakly increasing in $K$ and $B$; and
\item under the uniform-$\theta$ microfoundation of Section~\ref{subsec:two-model-demand} with a single quality tier, $A=\bar\theta\cdot Q$ and $B=\bar\theta\cdot Q/N$, so
\begin{equation}n^{*}=\left\lfloor\frac{K}{a}\cdot\sqrt{\frac{\bar\theta\cdot Q}{N\cdot(\rho+\Lambda)\cdot\mathcal F}}\right\rfloor,\end{equation}
which falls with the purchaser mass $N$ holding capacity fixed; whereas if capacity is slack, the standard Cournot free-entry count satisfies $n^{*}+1=\left\lfloor(A-b-a\cdot c)/\sqrt{B\cdot(\rho+\Lambda)\cdot\mathcal F}\right\rfloor$, which rises with $N$, proportionally to $\sqrt N$ (derived in Appendix~\ref{app:slack-capacity}).
\end{enumerate}
\end{proposition}

Parts~(1)--(3) follow from the strict monotonicity of $\pi(n)$: the entry condition $\pi(n)\geq(\rho+\Lambda)\cdot\mathcal F$ rearranges to $n\leq(K/a)\sqrt{B/((\rho+\Lambda)\mathcal F)}$, and the largest such integer is the unique fixed point at which every active laboratory covers its entry cost while one more entrant would not. Part~(4) substitutes the microfounded demand parameters into the binding-capacity formula and into the textbook free-entry condition $\pi^{\mathrm{slack}}(n)=(A-b-a\cdot c)^{2}/(B\cdot(n+1)^{2})\geq(\rho+\Lambda)\cdot\mathcal F$ for the slack case.

Part~(4) deserves emphasis because it inverts a standard intuition. The classic sunk-cost logic says the number of firms grows like the square root of market size over fixed cost, so bigger markets support more entrants. Here that holds only while capacity is slack. While capacity binds, total output is pinned at $K/a$, so a larger purchaser mass raises the price level and the scarcity value of compute---but the scarcity value flows into the wholesale price $r^{*}$ and accrues to hyperscalers, while the laboratories' Cournot markup $B\cdot K/(a\cdot n)$, the piece of the surplus that actually repays training costs, is proportional to the demand slope at the constrained quantity and \emph{shrinks} as the market grows. In short: while GPUs bind, market growth funds hyperscalers, not entry. This is the entry-margin counterpart of Proposition~\ref{prop:capacity-neutrality}: under binding capacity, adding laboratories does not lower consumer prices, and growing the market does not add laboratories.

The comparative statics in parts~(2) and~(3) are the formal version of a familiar industrial comparison. Whether the frontier-model market looks more like commercial aircraft---a worldwide market that supports roughly two full-line manufacturers---or like ordinary software, where fixed costs are trivial next to the addressable profit pool, is in this framing the question of where $(\rho+\Lambda)\cdot\mathcal F$ sits relative to the pool $B\cdot K^{2}/a^{2}$. Larger training runs, faster displacement hazards, or a smaller addressable profit pool all push toward fewer, larger laboratories.

Two qualifications keep the result honest. First, $\mathcal F$ is sunk at entry, so the free-entry count is an ex ante object: a market can transiently carry more than $n^{*}$ laboratories, because entrants who paid on expectations that did not pan out rationally keep operating as long as the flow conditions of Section~\ref{sec:trailing-labs} hold---the training cost is gone either way. The flow conditions there govern who stays; the stock condition here governs who should arrive. Second, the proposition imposes a common $\mathcal F$ at quality parity, parallel to the deliberate symmetry of Proposition~\ref{prop:capacity-neutrality}. With heterogeneous research efficiency, Section~\ref{subsec:entry-cost} makes the fixed cost laboratory-specific, and free entry admits laboratories in descending order of $\eta_\ell$ until the marginal entrant's value no longer covers its own $\mathcal F_\ell$.

\subsection{Segmentation changes how many laboratories can exist}
\label{subsec:segmented-entry}

Section~\ref{subsec:task-segmentation} read task segmentation as a demand-side force that redistributes rents and shares among the laboratories that exist. The entry margin makes it do something stronger: it changes how many laboratories can exist, because a laboratory that dominates a defensible segment needs to recover $\mathcal F$ only from that segment's profit pool rather than from a contested share of the whole market's.

\begin{proposition}[Entry under segment-level dominance]\label{prop:segmented-entry}
Divide the uniform-$\theta$ market of Proposition~\ref{prop:free-entry}(4) into $S$ symmetric segments, each containing purchaser mass $N/S$, so that each segment's inverse demand is $P_k=A-S\cdot B\cdot Y_k$. Assign each active laboratory to one segment as its sole intended provider. Laboratories are at exact quality parity---so quality cannot be what protects a niche here, in contrast to Proposition~\ref{prop:segment-wta}'s general mechanism---and instead every purchaser in segment $k$ incurs a \emph{recurring, per-output} friction $s\geq0$ (the recurring component of Section~\ref{subsec:task-segmentation}'s $s_{imk}$; the discussion following the proof explains why a one-time relationship cost cannot play this role) when patronizing any laboratory other than the one assigned to $k$, so a rival can attract a segment-$k$ purchaser only by pricing at least $s$ below the assigned laboratory. Let aggregate capacity bind under this configuration, and let compute clear at a common wholesale price $r^{*}=(A-b)/a-2BK/a^{2}$. Then, provided
\begin{equation}\label{eq:friction-threshold}s\;\geq\;s_{\min}\equiv\frac{B\cdot K}{a}\;=\;P_k^{*}-\left(b+a\cdot r^{*}\right)\end{equation}
(the friction must be at least the full segment markup, $P_k^{*}=A-B\cdot K/a$ being the segmented equilibrium price of part~(2); see the derivation below),
and provided $K<a(A-b-a\cdot c)/(2B)$ (capacity actually binds under the segmented configuration; see the discussion after the proof for why this is a stricter requirement than the head-on binding threshold):
\begin{enumerate}
\item each laboratory's flow profit is
\begin{equation}\pi^{\mathrm{seg}}(S)=\frac{B\cdot K^{2}}{a^{2}\cdot S},\end{equation}
which is $S$ times what each of the same $S$ laboratories would earn competing head-on for the whole market;
\item the purchaser-facing price is $A-B\cdot K/a$ in both configurations, so the capacity-binding neutrality of Proposition~\ref{prop:capacity-neutrality} extends: segmentation changes only the division of rents, lowering the wholesale compute price and moving fixed-capacity rent from hyperscalers to laboratories; and
\item free entry supports any $S$ with $\pi^{\mathrm{seg}}(S)/(\rho+\Lambda)\geq\mathcal F$, i.e., up to
\begin{equation}\bar S=\left\lfloor\frac{B\cdot K^{2}}{a^{2}\cdot(\rho+\Lambda)\cdot\mathcal F}\right\rfloor\end{equation}
laboratories, which is, up to integer parts, the \emph{square} of the head-on count $n^{*}$ of Proposition~\ref{prop:free-entry}. As with Proposition~\ref{prop:free-entry}, $\pi^{\mathrm{seg}}(S)$ is the competitive-upstream flow profit; under bargained upstream conduct (Appendix~\ref{app:bargaining}) it is correspondingly larger and $\bar S$ is a lower bound ($\bar S$ is a viability ceiling, not an equilibrium count; see the discussion following the proof); and
\item the friction condition~\eqref{eq:friction-threshold} is what protects a segment against \emph{already-trained} rivals---in particular the other $S-1$ laboratories, whose training cost is sunk and who, because compute is fungible across segments, can redeploy their models into segment $k$ at compute opportunity cost $r^{*}$ alone. Against a laboratory that has \emph{not yet} trained a frontier model, the operative barrier is $\mathcal F$ itself: a quality-parity entrant that pays $\mathcal F$ and contests one segment earns post-entry flow profit (for $s\leq B\cdot K/a$; above that it attracts no purchasers and earns nothing)
\begin{equation}\label{eq:entrant-duopoly-profit}\pi^{E}(S,s)=\frac{4\cdot S\cdot\left(B\cdot K/a-s\right)^{2}}{B\cdot(3S+1)^{2}}\;\leq\;\pi^{E}(S,0)=\frac{4\cdot S^{2}}{(3S+1)^{2}}\cdot\pi^{\mathrm{seg}}(S)\;<\;\frac{4}{9}\cdot\pi^{\mathrm{seg}}(S),\end{equation}
so fresh entry into an occupied segment is strictly unprofitable---at \emph{any} friction level, including $s=0$---whenever $(\rho+\Lambda)\cdot\mathcal F>\pi^{E}(S,0)$. Because $\pi^{E}(S,0)<\pi^{\mathrm{seg}}(S)$, the range of fixed costs at which the $S$ incumbents are viable (part~3) while every fresh entrant is deterred with no friction assumption at all is nonempty for every $S$.
\end{enumerate}
\end{proposition}

The profit and price derivation is three lines, and is unaffected by how the niche is protected: a segment monopolist facing marginal cost $b+a\cdot r$ sets $A-2\cdot S\cdot B\cdot y_k=b+a\cdot r$; the common wholesale price adjusts until the symmetric allocations exhaust capacity, $y_k=K/(a\cdot S)$; the markup is then $S\cdot B\cdot y_k=B\cdot K/a$ regardless of $S$, and flow profit is markup times output, giving part~(1). The segment price $A-S\cdot B\cdot K/(a\cdot S)=A-B\cdot K/a$ gives part~(2); the implied wholesale price $r^{*}=(A-b)/a-2\cdot B\cdot K/a^{2}$ is below its head-on counterpart from Appendix~\ref{app:cournot} for any $n\geq2$, which is the rent transfer. Part~(3) substitutes part~(1) into the entry condition and compares with Proposition~\ref{prop:free-entry}(1).

The protection condition~\eqref{eq:friction-threshold} needs its own derivation, and earlier drafts of this proposition instead invoked ``the protected-niche configuration of Proposition~\ref{prop:segment-wta}'' directly---which is not available here. Under this proposition's own symmetric-quality assumption, a rival's value schedule in segment $k$ is identical to the assigned laboratory's, so absent any friction a rival's marginal value of a first unit in segment $k$ would equal the incumbent's and strictly exceed $r^*$, \emph{violating} rather than satisfying Proposition~\ref{prop:segment-wta}'s protection condition $v_{\ell k}(0)\leq r^*$. Quality parity is therefore not compatible with protected niches on its own; the model needs the friction to do that work, and Section~\ref{subsec:task-segmentation} introduces $s_{imk}$ without ever pinning down a magnitude it must reach. Here we can. By the utility specification $u_{imk}=\theta_{ik}Q_{mk}-P_{mk}-s_{imk}$ with $s_{imk}=s$ for every non-assigned laboratory in segment $k$ and $s_{imk}=0$ for the assigned one, a rival can attract a purchaser only by pricing at $P_k^*-s$ or below, where $P_k^*=A-B\cdot K/a$ is the segmented equilibrium price from part~(2). The rival's gross marginal value per FLOP of a first unit in segment $k$---the object $v_{\ell k}(0)$ of Proposition~\ref{prop:segment-wta}, revenue net of non-compute cost per FLOP---is then
\begin{equation}v_{\mathrm{rival},k}(0)=\frac{P_k^*-s-b}{a},\end{equation}
and Proposition~\ref{prop:segment-wta}'s protection condition $v_{\mathrm{rival},k}(0)\leq r^*$ rearranges to $s\geq P_k^*-b-a\cdot r^*=B\cdot K/a=s_{\min}$, which is condition~\eqref{eq:friction-threshold}. (Watch the units here: comparing the rival's per-output margin \emph{net} of compute cost, $(P_k^*-s)-\kappa$ with $\kappa=b+a\cdot r^*$, against the per-FLOP price $r^*$ would double-count the compute cost and mix per-output with per-FLOP quantities, yielding the laxer but incorrect threshold $B\cdot K/a-r^*$; the contested-segment equilibrium computed for part~(4) below confirms the correct threshold independently, since a rival's post-entry profit $\pi^{E}(S,s)$ is strictly positive for every $s<B\cdot K/a$ and reaches zero exactly at $s=B\cdot K/a$.) The threshold has a sharp economic reading: $B\cdot K/a$ is exactly the segment markup of part~(1), so the friction must absorb the \emph{entire} margin the niche sustains. This is ordinary contestability logic. A rival at quality parity produces at the same compute-inclusive cost $\kappa=b+a\cdot r^*$ as the incumbent, so any price further than $s$ above $\kappa$ is profitably undercuttable; the incumbent's price exceeds $\kappa$ by precisely $B\cdot K/a$. It follows that under quality parity the segmented configuration is \emph{never} self-protecting: $s_{\min}>0$ whenever the markup is positive, the required friction scales one-for-one with the profitability of the niche it defends, and condition~\eqref{eq:friction-threshold} is the minimum recurring switching, latency, privacy, or contractual cost that must separate purchasers from an equal-quality rival for sole provision to be consistent with Proposition~\ref{prop:segment-wta}---not an assumption asserted at an unstated magnitude.

Note also what this derivation does and does not assume about who the rival is. It prices the rival's first unit at the compute opportunity cost $r^{*}$ and assumes quality parity---nothing else. It therefore applies verbatim to the other $S-1$ assigned laboratories: their training cost is sunk, compute is fungible across segments (Section~\ref{subsec:task-segmentation}), and cross-segment separability means an incursion into segment $k$ leaves their home-segment profits untouched, so the only question an already-trained rival faces is exactly the marginal redeployment question the condition prices, and $\mathcal F$ appears nowhere in it. Fixed training costs provide no protection whatsoever against these rivals; the friction is the entire defense. Symmetrically, the derivation says nothing about a laboratory with no trained model, which cannot supply a first unit in segment $k$ at any friction level without first paying $\mathcal F$; part~(4) prices that threat instead, and there the roles reverse: the fixed cost is the entire defense, and the friction is merely helpful. Two things remain genuinely unmodeled and are not claimed here: which laboratory is assigned to which segment (we take the assignment, like the number and size of segments, as a primitive of market structure rather than the outcome of an entry-into-segments game), and any account of why $s_{imk}$ would actually satisfy~\eqref{eq:friction-threshold} for a real market rather than falling short of it.

Part~(4) is derived in the same conduct regime as everything else in this section: Cournot within the contested segment, price-taking in $r$, capacity binding. Let a quality-parity entrant contest segment $1$ while the other $S-1$ segments remain monopolized, and write $\tilde D\equiv A-b-a\cdot r$. Each unchallenged monopolist supplies $y_k=\tilde D/(2SB)$; in segment $1$ the incumbent's and entrant's first-order conditions are $2SBy_I+SBy_E=\tilde D$ and $SBy_I+2SBy_E=\tilde D-s$ (the entrant sells at $P_1-s$), giving $y_I=(\tilde D+s)/(3SB)$ and $y_E=(\tilde D-2s)/(3SB)$. Capacity clearing, $a\left[(S-1)\tilde D/(2SB)+(2\tilde D-s)/(3SB)\right]=K$, pins down $\tilde D=(6SBK/a+2s)/(3S+1)$; the entrant's markup is $SBy_E=2S(BK/a-s)/(3S+1)$, and its flow profit $SBy_E^{2}$ is equation~\eqref{eq:entrant-duopoly-profit}. Three features of this equilibrium are worth recording. First, entry \emph{raises} the wholesale price: $\tilde D<A-b-a\cdot r^{*}$ for every $s<BK/a$, so the post-entry clearing price exceeds the segmented $r^{*}$---an additional competitor weakens output restriction, and with $K$ fixed the extra compute demand accrues to hyperscalers, making entry partly self-defeating for the same reason market growth fails to fund entry in Proposition~\ref{prop:free-entry}(4). Second, capacity still binds post-entry on the proposition's whole domain: $\tilde D\leq2BK/a<A-b-a\cdot c$ for every $s\in[0,BK/a]$ given $K<a(A-b-ac)/(2B)$, so the comparison never leaves the binding regime. Third, $\pi^{E}(S,s)$ is strictly decreasing in $s$ and vanishes exactly at $s=s_{\min}=BK/a$, which verifies the protection threshold from the equilibrium side: below $s_{\min}$ a contested-segment equilibrium exists and the segmented configuration is not an equilibrium (the undercutting deviation priced above is profitable); at or above $s_{\min}$ the contested equilibrium disappears ($y_E\leq0$) and the segmented configuration is self-enforcing. The deviation condition and the equilibrium condition are exact complements, with no window of multiplicity between them.

The friction in condition~\eqref{eq:friction-threshold} must also be the right \emph{kind} of friction, which is why the proposition's statement insists on the recurring component of $s_{imk}$. The derivation above subtracts $s$ from the rival's price on every unit, so it requires a per-output cost: latency, privacy costs, per-task compliance or contractual overhead. A one-time relationship cost is a different primitive---the evaluation cost $E_{ij}$ of Appendix~\ref{app:subscriptions}, paid once to add a provider to a subscription portfolio, after which routing among subscribed providers is frictionless---and it cannot stand in for $s$ here, for two reasons. First, it amortizes: to a purchaser expecting a discounted volume of $T$ outputs from the segment, a one-time cost $E$ is a per-output wedge of only $E/T$, so for high-volume purchasers---precisely the purchasers that matter most for compute demand---even a large one-time cost falls far short of a per-output $s_{\min}$ equal to the full markup. Second, it is extinguished by multi-homing: a purchaser whose task mix spans segments pays $E_{ij}$ to add a second provider whenever the incremental-value condition of Appendix~\ref{app:subscriptions} clears, and from then on faces no wedge at all against routing segment-$k$ tasks to that provider. The proposition's symmetric partition of purchasers into $S$ disjoint segments quietly assumes this case away: every purchaser has tasks in one segment only, so no purchaser has an Appendix-\ref{app:subscriptions} reason to multi-home, and the persistent-multi-homing prediction of Section~\ref{subsec:task-segmentation} is switched off by construction. Stated honestly, then: segment protection here is a claim about single-homed purchasers facing recurring frictions. For purchasers who do multi-home, the one-time component of their frictions is already sunk, only the recurring component continues to protect the niche, and protection erodes as multi-homing spreads---the more the market looks like Section~\ref{subsec:task-segmentation}'s portfolio picture, the less of the purchaser base condition~\eqref{eq:friction-threshold} actually covers. Subscription pricing itself sits between the two cases: the per-period fee $\phi_j$ of Appendix~\ref{app:subscriptions} is a recurring cost of \emph{maintaining} access to a second provider, so it does act as a genuine recurring wedge---but at the relationship level, amortizing over the purchaser's volume like $\phi_j/T$, and so protecting niches against low-volume purchasers while doing little against high-volume ones.

Both capacity-binding thresholds matter because the two configurations do not bind at the same $K$, and the proposition's domain restriction above is doing real work. Without a capacity constraint, a segment monopolist's unconstrained output solves $A-2SBy_k=b+ac$, so aggregate desired compute is $a\cdot S\cdot y_k=a(A-b-ac)/(2B)$; the segmented configuration binds only below this level. Head-on Cournot among the same $S=n$ laboratories binds at the laxer threshold $K<aS(A-b-ac)/(B(S+1))$ (equation~\eqref{eq:cournot-slack-threshold}, Appendix~\ref{app:cournot}), which exceeds the segmented threshold for every $S\geq1$ (with equality only at $S=1$): segment monopolists individually restrict output further than $S$ Cournot rivals sharing one market, so less capacity is needed to make them bind. For $K$ strictly between these two thresholds, head-on Cournot capacity binds but the segmented configuration would be slack---$r^*=c$ there, parts~(1)--(2) both fail as stated, and comparing $\pi^{\mathrm{seg}}$ against $\pi(n)$ no longer compares two binding-capacity configurations. The proposition is stated, and should be read, only for $K$ below the segmented threshold; the excluded window between the two thresholds is a real gap in the comparison rather than a technicality, and we do not attempt to characterize outcomes there. The fully slack region above both thresholds, by contrast, is characterized in Appendix~\ref{app:slack-capacity}: there the protection threshold~\eqref{eq:friction-threshold} rises to the unconstrained monopoly markup and segmentation stops being price-neutral.

Two forces explain why the head-on count misses what segmentation adds, on the domain where the proposition holds. Differentiation protects markups: a segment leader prices as a monopolist against its own residual demand rather than as one of $n$ Cournot rivals on common demand, which is the Lerner logic of Section~\ref{sec:oligopoly} operating segment by segment. And differentiation partitions the cost-recovery problem: the fixed cost is compared against a whole segment's pool, not against a $1/n$ share of an eroding aggregate pool. The exact square relationship in part~(3) is special to the symmetric linear case and should not be quoted as more than an illustration, but the direction is general: fragmentation of demand into defensible niches raises the number of laboratories that a given profit pool can carry past the fixed-cost bar. Segment differentiation therefore does not merely redistribute rents among incumbents, as in Section~\ref{subsec:task-segmentation}; it relaxes the natural-oligopoly bound itself, on segments where a recurring friction at least as large as the segment markup actually separates single-homed purchasers from equal-quality rivals.

There is a countervailing force, and the friction condition changes where it operates---but does not eliminate it. The obvious version of the concern is marginal invasion: $\mathcal F$ is paid once per laboratory, and compute fungibility means an incumbent that has already sunk its training cost could serve an adjacent segment at no incremental training cost. Under condition~\eqref{eq:friction-threshold} that is exactly what part~(4)'s first clause blocks: no incremental training cost is needed, but the friction makes the redeployment unprofitable at the going prices, so the assignment, once in place, is stable against it. The force survives at a stage the model does not represent: the formation of the assignment itself. Nothing in the model explains why $S$ laboratories each pay $\mathcal F$ rather than one laboratory paying $\mathcal F$ once and becoming the assigned---friction-free---provider of every segment. A single laboratory assigned to all $S$ segments earns $S\cdot\pi^{\mathrm{seg}}(S)=B\cdot K^{2}/a^{2}$ per period, the head-on monopoly profit $\pi(1)$ of Proposition~\ref{prop:free-entry}, for one training cost; it is protected against fresh entrants by exactly the same condition as in part~(4), since price-taking in $r^{*}$ makes its first-order conditions segment-separable and the contested-segment equilibrium is unchanged; and no purchaser ever bears the friction on the equilibrium path. The spanning configuration is therefore strictly privately better for the laboratory that achieves it and no less stable once achieved. Nothing in the model's primitives breaks the tie the other way: quality is at parity by assumption, the task comparative advantage $\chi_{\ell,k}$ of Section~\ref{sec:persistent-advantage} is an endowment rather than a per-segment investment, and $s_{imk}$ is a primitive of purchaser utility attached to an assignment the model takes as given. What would break it is a per-segment cost of \emph{becoming} a segment's friction-free provider---segment-specific integration, distribution, sales, or fine-tuning investment, that is, an account in which the assignment, or the advantage $\chi_{\ell,k}$, must be bought segment by segment---and the model contains no such choice variable. Part~(3)'s $\bar S$ must accordingly be read as a viability ceiling, not an equilibrium count: it is the largest number of laboratories the segmented profit pool can carry past the fixed-cost bar, and it is attained only if the unmodeled assignment process actually spreads segments across distinct laboratories. The model shows that every configuration from one laboratory spanning all segments to $\bar S$ laboratories holding one each is stable and viable under the stated conditions; it does not select among them, and we do not claim otherwise.

\subsection{What is derived and what is conjectured}
\label{subsec:entry-honesty}

The free-entry fixed point, the fourfold natural-monopoly window, and the binding-versus-slack reversal of the market-size comparative static are all derived (the slack side of the reversal is worked out in Appendix~\ref{app:slack-capacity}), and derived from machinery already in the paper: the quality laws of Section~\ref{sec:technology-regimes} supply the compute cost of entry, the training-compute condition of Section~\ref{sec:model-race} prices it, and the Cournot expressions of Appendix~\ref{app:cournot} supply post-entry profits. The segmentation result of Proposition~\ref{prop:segmented-entry} is derived from the same machinery, but---unlike the other three---only \emph{conditional on a configuration that the proposition now states explicitly rather than presuming}: quality parity plus a recurring, per-output friction $s$ at or above the full segment markup $B\cdot K/a$ (equation~\eqref{eq:friction-threshold}), and aggregate capacity binding below the segmented-specific threshold $a(A-b-ac)/(2B)$, which is stricter than the head-on binding threshold. Both conditions are necessary for the configuration to be internally consistent with Proposition~\ref{prop:segment-wta}; neither is guaranteed to hold for any particular real market, and the proposition does not claim it does---it says what follows \emph{if} they hold, and gives the threshold values so the ``if'' is checkable rather than asserted. This is a materially different, and more honest, status than ``derived'' full stop: the entry-count and natural-monopoly results above are unconditional consequences of the model's other primitives, while the segmentation result is conditional on a friction and a capacity regime whose magnitudes the paper does not estimate. Three further qualifications, worked out in the discussion after Proposition~\ref{prop:segmented-entry}, bound what even the conditional result covers. First, the friction condition and the fixed cost protect against different rivals, and the division of labor is now derived rather than blurred: $s\geq B\cdot K/a$ is what deters the other already-trained laboratories, for whom $\mathcal F$ is sunk and irrelevant, while a not-yet-trained entrant is deterred by $\mathcal F$ alone---no friction needed---whenever $(\rho+\Lambda)\cdot\mathcal F$ exceeds the contested-segment profit $\pi^{E}(S,0)$ of Proposition~\ref{prop:segmented-entry}(4), a bound that sits below the viability bar for every $S$. Second, the friction that clears the threshold must be recurring and the purchasers single-homed: a one-time relationship cost (the $E_{ij}$ of Appendix~\ref{app:subscriptions}) amortizes over usage volume and is already sunk for any purchaser who multi-homes, so segment protection is a claim about single-homed purchasers facing per-output frictions, and it erodes as multi-homing spreads. Third---the largest remaining gap---the proposition establishes that the $S$-laboratory assignment is self-enforcing and viable, not that it is selected: one laboratory paying $\mathcal F$ once and holding the friction-free assignment in every segment earns $S\cdot\pi^{\mathrm{seg}}(S)$ under the same protection, so $\bar S$ is a ceiling on how many laboratories segmentation can carry, and the model contains no mechanism---such as a per-segment cost of acquiring a segment's assigned-provider position---that predicts the ceiling is reached rather than one laboratory spanning every segment. Beyond this, two further things are deliberately not claimed. First, no numerical prediction of the number of laboratories: equation~\eqref{eq:entry-cost} prices only the compute component of the fixed cost, and the symmetric linear demand system is an illustration, not a calibration. Second, the entry quality $q$ is held fixed above, while in reality it is chosen competitively, and the choice escalates: each entrant's incentive to outspend rivals on training raises the quality bar, and with it $\mathcal F$, in step with the profit pool. This is precisely the mechanism by which \citet{sutton1991} argues that growing markets need not fragment---sunk outlays rise endogenously with market size---and the steep rise of $\Gamma_\ell(q)$ in the entry quality, divergent near the asymptotes and convex for $\beta<1$ under continuing improvement (Section~\ref{subsec:entry-cost}), is exactly the channel through which such escalation would operate here. We do not solve the entry-quality game, so that implication remains qualitative.

\section{Implications for Market Structure}
\label{sec:implications}

The model predicts a layered and dynamic market structure.

\begin{enumerate}
\item One laboratory or a small set of frontier laboratories serves the highest-value demand and can bid aggressively for scarce compute.
\item Near-frontier laboratories remain viable by charging lower prices, achieving greater productivity on some tasks, serving differentiated market segments, or becoming secondary subscription providers.
\item Far-behind laboratories cannot generally compete for scarce compute in the undifferentiated inference market. They must specialize, rely on subsidized or previously contracted hardware, provide services around commoditized models, or exit frontier-model development.
\item Hyperscalers earn scarcity rents when capacity is binding, but their share depends on competition among laboratories for compute and on the laboratories' ability to exercise oligopsony.
\item Oligopolistic laboratories can retain positive downstream markups even when compute capacity is abundant, so greater GPU supply is not sufficient by itself to guarantee low consumer prices.
\item Continuing frontier progress generates temporary leadership, lumpy market-share shifts, and repeated cycles in which model-provider rents rise after breakthroughs and hyperscaler bargaining power rises during catch-up.
\end{enumerate}

The model's central relationships are:

\begin{enumerate}[label=\textbf{\arabic*.}]
\item Compute is allocated by marginal private value per constrained FLOP.
\item Consumer price equals physical cost plus GPU scarcity rent plus model-provider markup.
\item Laboratory rent depends on differential productivity, downstream market power, and upstream bargaining power.
\item Hyperscaler rent depends on capacity scarcity and competition among laboratories for compute.
\end{enumerate}

A persistent productivity leader is likely to receive a persistent share of inference compute. Complete dominance requires more: the leader's marginal use of the final available FLOP must remain more valuable than every follower's best initial use. Under oligopoly, even incomplete dominance can sustain substantial model-provider rents, and convergence in technical performance need not imply convergence to marginal-cost pricing.

\section{Vertical Integration and Heterogeneous Hyperscalers}
\label{sec:vertical-integration}

The baseline treats laboratories and hyperscalers as separate firms meeting in a market for compute. In practice the two layers are frequently integrated, and this is worth taking seriously because it is the model's largest realism gap. It is, however, less damaging to the results than it first appears, because integration bears very differently on the allocation results than on the division-of-rents results.

The allocation rule and the shadow price survive integration. The condition $v_m=c+\lambda$ and the capacity shadow price $\lambda$ depend only on marginal net revenue per FLOP and the physical constraint, not on whether a market transfer actually changes hands. An integrated firm internalizes the wholesale price $r$ as an internal opportunity cost of compute and allocates its own FLOPs to their highest-value use in exactly the same way; the boxed allocation condition of Section~\ref{sec:allocation} is unchanged. What integration changes is the division of rents. When a laboratory and a hyperscaler are the same firm there is no arms-length price at which one pays the other, so the split between the hyperscaler scarcity component $a_\ell\cdot(r_\ell-c)$ and the laboratory markup $\mu_\ell$ ceases to be a market outcome and becomes an internal accounting choice. The bargaining and auction results of Appendix~\ref{app:bargaining}, which describe how the match-specific surplus $\widehat v_\ell-\bar v_h$ is split, then apply only to pairs that remain at arms length; an integrated pair simply retains the whole surplus. The empirical counterpart is that transfer prices for compute booked inside integrated firms need not reveal the shadow value of capacity, and that measured rent shares across the two layers are partly an artifact of where corporate boundaries happen to fall.

A natural reduced-form extension replaces the representative hyperscaler with heterogeneous suppliers $h$, each characterized by capacity $K_h$, operating cost $c_h$, and an outside- or internal-demand parameter $o_h$. This captures neoclouds, large diversified clouds, captive internal demand, and cost differences across hardware vintages without modeling the semiconductor supply chain: the single clearing price $r^*$ is then replaced by a supply schedule traced out by the ordered marginal costs and outside options of the individual suppliers, and vertical integration appears as a supplier whose $o_h$ reflects captive demand from an affiliated laboratory rather than an external market price.

\section{Empirical Implications}
\label{sec:empirical}

The model is stylized, but it generates several predictions that are, at least in principle, falsifiable against existing and emerging data on the LLM market. We state each as a testable claim tied to the specific model object that produces it.

\paragraph{1. Under binding capacity, prices track compute supply, not lab entry.} When capacity binds, the equilibrium consumer price depends on aggregate compute $K$ but is invariant to the number of laboratories serving the market, holding the frontier fixed (the capacity-binding neutrality result of Proposition~\ref{prop:capacity-neutrality} in Section~\ref{sec:oligopoly}, proved in Appendix~\ref{app:cournot}, where $P^*=A-B\cdot K/a$ contains no $n$). The observable implication is that, in periods when inference demand exceeds capacity at near-cost prices, the price for a given quality tier should respond to changes in aggregate GPU supply while being roughly insensitive to additional providers entering. This can be checked against the price, performance, and supply series assembled by \citet{demirerEtAl2025}.

\paragraph{2. The leader's markup compresses on competitor releases.} The leader's markup satisfies the Lerner condition $(P_L-\kappa_L)/P_L=1/\epsilon_{LL}$ of Section~\ref{sec:oligopoly}. A close-substitute release by a rival raises the leader's own-price elasticity $\epsilon_{LL}$ and compresses its markup. The prediction is that a frontier model's price, relative to its quality, should fall discretely in the period following a rival's near-frontier release rather than declining on a smooth schedule. This is checkable against the release-timed price-cut patterns documented by \citet{demirerEtAl2025}.

\paragraph{3. For trailing labs, subscription share $\gg$ usage share $\gg$ compute share.} Because multi-homing lets a provider hold a broad subscriber base while drawing only a small fraction of inference traffic, a trailing laboratory's subscription (subscriber) share should exceed its usage (traffic) share, which should in turn exceed its compute share (Appendix~\ref{app:subscriptions}). This ordering is checkable against the multi-homing and post-release adoption evidence in \citet{fradkin2025demand}.

\paragraph{4. Hyperscaler rental prices rise after major releases.} A major release shifts the leading laboratory's private marginal value $\widehat v_L$ upward, and once competitors catch up it raises the best losing bidder's value $\widehat v_F$, which sets the clearing compute price $r^*\simeq\max\{c,\widehat v_F\}$. The same force appears in the hyperscaler capital-expenditure Euler condition of Section~\ref{sec:model-race}, where higher expected future scarcity rents $\lambda_{h,t+s}$ raise current investment. The prediction is that wholesale rental prices for inference-capable compute, and hyperscaler capital expenditure, should rise following breakthrough releases that expand inference demand. The industry-structure forces behind this concentration of value in compute are discussed by \citet{korinekVipra2025}.

\paragraph{5. Crossing into slack capacity reverses which variable moves prices.} Predictions~1 and~4 are conditional on binding capacity. Appendix~\ref{app:slack-capacity} derives the mirror predictions for the slack regime: the wholesale margin $r^{*}-c$ collapses, so rental prices for inference-capable compute fall to physical operating cost and stop contributing to capital recovery; measured utilization falls below capacity; and the consumer price of a given quality tier begins responding to laboratory entry and near-frontier releases---the ordinary Cournot dependence on $n$---while becoming insensitive to aggregate GPU supply. Because the wholesale margin is itself the observable regime marker, the joint prediction is sharper than either regime's alone: a market crossing the boundary should switch which right-hand-side variable consumer prices respond to, from capacity growth to entry, at the same time as the wholesale margin and utilization fall. The price and supply series of \citet{demirerEtAl2025} are in principle suited to detecting the switch.

\section{Limitations}
\label{sec:limitations}

Several mechanisms that matter for the LLM market are outside the model.

\paragraph{A patchwork of conduct assumptions, made explicit rather than resolved.} The paper does not use one conduct model throughout. Downstream, laboratories are Cournot quantity-setters where capacity binds (Sections~\ref{sec:oligopoly}--\ref{sec:natural-oligopoly}) but unconstrained Bertrand price-setters where the markup and elasticity dynamics of Proposition~\ref{prop:breakthrough} are derived (Section~\ref{subsec:price-competition} states that scope restriction at the point Bertrand is first used, and Appendix~\ref{app:two-model} adds a capacity-constrained Cournot derivation for the results that need one). Upstream, laboratories are price takers in the compute market in Propositions~\ref{prop:capacity-neutrality}, \ref{prop:free-entry}, and \ref{prop:segmented-entry}, while Appendix~\ref{app:bargaining} argues that a market with few laboratories and few hyperscalers is realistically a bilateral oligopoly with bargaining and oligopsony power on both sides. Each proposition states which conduct assumption it uses at the point it is first invoked, and Section~\ref{sec:oligopoly} and Appendix~\ref{app:bargaining} cross-reference each other on how the results would change under the other assumption, but the paper does not pick a single conduct model and derive everything from it; a reader should treat each numbered result as conditional on the conduct assumption stated where it appears, not as a description of the industry's actual conduct. A parallel patchwork runs across the capacity regime itself: most numbered results are stated for binding capacity, and Appendix~\ref{app:slack-capacity} collects their slack-capacity counterparts in one place rather than unifying the two regimes into single statements.

\paragraph{No labor-market side.} Downstream willingness to pay is a primitive here, and we do not model the labor market that generates much of it. As prices rise, firms will substitute between human and model labor, and in some cases humans may remain cheaper than models; this is noted only in a footnote in the introduction and deserves separate treatment.

\paragraph{No open-weight or self-hosted models.} We model all inference as purchased from laboratories, but open-weight models that purchasers can self-host act as a competitive fringe. On commodity tasks they cap both the wholesale compute price $r^*$ and the laboratory markup $\mu_\ell$ at the cost of self-hosting, effectively turning the trailing-laboratory tier into a zero-price competitor for undifferentiated work. This is a genuinely missing mechanism rather than a harmless normalization, and it would blunt several of the rent results on commoditized tasks.

\paragraph{Exogenous demand growth.} Demand is taken as given at each date. In reality adoption is endogenous to price, capability, and complementary investment, and the feedback from falling quality-adjusted prices back to demand is a first-order driver of the very scarcity the model studies.

\paragraph{No agentic or inference-time-compute distinction.} We treat FLOPs per output $a_m$ as a technological primitive that falls with productivity. Inference-time-compute and agentic workflows instead raise the compute required per completed task, which could offset or reverse the downward trend in $a_m$ and change the allocation over time.

\paragraph{A single homogeneous FLOP, and no separate power constraint.} Compute is a scalar constrained resource. Appendix~\ref{app:extensions} partially relaxes this by allowing multiple constrained resources, but the baseline abstracts from heterogeneous accelerators, memory bandwidth, and interconnect. Relatedly, we impose no energy or power constraint separate from FLOPs, though power availability and grid interconnection may bind before compute does.

\section{Conclusion}
\label{sec:conclusion}

A fixed stock of inference hardware creates an upstream scarcity market in which models compete for compute according to the commercial value they generate per constrained FLOP. Persistent laboratory advantage should be represented multiplicatively because a fixed additive increment becomes negligible relative to an exponentially advancing frontier. This advantage shifts a laboratory's marginal-value schedule upward, but does not by itself imply complete market dominance.

The integrated-industry benchmark identifies the joint-profit-maximizing allocation of compute, while the oligopoly model shows how downstream market power changes prices and rent division. This joint-profit benchmark coincides with the welfare-maximizing allocation only under perfect price discrimination or price-taking behavior; in general it equalizes marginal net revenue rather than marginal social value per FLOP. Purchaser-facing prices contain physical cost, a possible hyperscaler scarcity rent, and a model-provider markup. When capacity is tightly binding, oligopoly may primarily transfer rents from hyperscalers to laboratories. When capacity is slack, hyperscaler scarcity rents disappear but laboratory markups can persist.

The dynamic model connects current inference, training investment, and future GPU capacity. Lumpy model releases create temporary leadership rents and induce both laboratories and hyperscalers to invest. Under a common global asymptote, differential model rents tend to disappear, though oligopoly markups may remain. Firm-specific asymptotes preserve durable differentiation. Continuing improvement creates repeated episodes of temporary leadership, with laboratories capturing more surplus after breakthroughs and hyperscalers gaining bargaining power as competitors catch up.

\appendix

\section{Bargaining, Auctions, and Oligopsony}
\label{app:bargaining}

The main text uses a competitive upstream auction as a benchmark: Propositions~\ref{prop:capacity-neutrality}, \ref{prop:free-entry}, and \ref{prop:segmented-entry} all treat laboratories as price takers with respect to $r^*$, which each proposition now states at the point it is first used. This appendix gives a bilateral Nash-bargaining formulation and clarifies the effect of concentrated laboratory demand for compute; the pointers at those three propositions send the reader here for how the competitive-benchmark results change once laboratories have bargaining power.

Consider a marginal block of compute allocated by hyperscaler $h$ to laboratory $\ell$. Let $\widehat v_\ell$ be the laboratory's private marginal profit per FLOP. Let $\bar v_h$ denote the hyperscaler's value from its best alternative use of the same FLOP, including physical operation or sale to another laboratory:

\begin{equation}\bar v_h\equiv\max\{c,\widehat v_F\}.\end{equation}

The match-specific surplus is

\begin{equation}S_{\ell h}=\widehat v_\ell-\bar v_h.\end{equation}

Suppose the laboratory has Nash bargaining weight $\zeta\in[0,1]$. If $r_{\ell h}$ is the negotiated wholesale price per FLOP, the Nash solution maximizes

\begin{equation}\left(\widehat v_\ell-r_{\ell h}\right)^\zeta\cdot\left(r_{\ell h}-\bar v_h\right)^{1-\zeta}.\end{equation}

The negotiated price is

\begin{equation}r_{\ell h}=\zeta\cdot\bar v_h+\left(1-\zeta\right)\cdot\widehat v_\ell.\end{equation}

The laboratory receives

\begin{equation}\pi_\ell^F=\widehat v_\ell-r_{\ell h}=\zeta\cdot\left(\widehat v_\ell-\bar v_h\right),\end{equation}

while the hyperscaler's incremental gain over its outside option is

\begin{equation}r_{\ell h}-\bar v_h=\left(1-\zeta\right)\cdot\left(\widehat v_\ell-\bar v_h\right).\end{equation}

The hyperscaler's total operating margin over physical cost is

\begin{equation}r_{\ell h}-c=\underbrace{\bar v_h-c}_{\text{general scarcity rent}}+\underbrace{\left(1-\zeta\right)\cdot\left(\widehat v_\ell-\bar v_h\right)}_{\text{share of match-specific surplus}}.\end{equation}

With many substitutable hyperscalers competing for a leading laboratory, $\zeta$ may be large and the laboratory may capture most match-specific surplus. With fixed aggregate hardware and many laboratories bidding for it, $\bar v_h$ may be close to the value generated by the best losing bidder, allowing the hyperscaler to capture most general scarcity rent.

If only a few laboratories purchase frontier compute, they may exercise oligopsony. They can reduce the wholesale compute price below the competitive-auction benchmark by restricting or coordinating demand, subject to the hyperscaler's outside options. A market with few laboratories and few hyperscalers is therefore a bilateral oligopoly: downstream oligopoly raises purchaser prices, while upstream oligopsony transfers part of the compute scarcity rent from hyperscalers to laboratories.

\section{Dynamic Operation by Trailing Laboratories}
\label{app:dynamic-loss}

Let $s_{jt}$ summarize trailing laboratory $j$'s strategic state, including model capability, customer relationships, deployment knowledge, data, reputation, and future access to capital or compute. Let inference activity be $x_{jt}$, and suppose the state evolves according to

\begin{equation}s_{j,t+1}=H(s_{jt},x_{jt})+\varepsilon_{j,t+1}.\end{equation}

The laboratory's value function is

\begin{equation}V_j(s)=\max_{x\geq0}\left\{\pi_j^{\mathrm{current}}(s,x)+\delta\cdot\mathbb E\left[V_j\left(H(s,x)+\varepsilon'\right)\right]\right\}.\end{equation}

Operating at level $x>0$ is preferable to non-operation when

\begin{equation}\pi_j^{\mathrm{current}}(s,x)-\pi_j^{\mathrm{current}}(s,0)+\delta\cdot\left[\mathbb E[V_j(s')\mid s,x]-\mathbb E[V_j(s')\mid s,0]\right]\geq0.\end{equation}

If current incremental profit is negative, operation is rational only when the discounted improvement in continuation value is large enough to offset the loss. The transition effect may reflect customer retention, useful deployment data, operational learning, reputation, distribution, or an increased probability of financing or successfully training a future model. These benefits should not be assumed. If inference deployment contributes little to frontier research capability, directing the same resources toward research may dominate subsidized inference operation.

\section{Subscription Pricing and Multi-Homing}
\label{app:subscriptions}

Provider $m$ offers a subscription contract $(\phi_m,\rho_m)$, where $\phi_m\geq0$ is a fixed subscription fee and $\rho_m$ is the marginal price of a standardized output after subscribing. A non-subscriber may instead pay spot price $P_m^{\mathrm{spot}}$, with

\begin{equation}\rho_m<P_m^{\mathrm{spot}}.\end{equation}

The fixed fee gives the provider guaranteed revenue in exchange for discounted future usage. If purchaser $i$ expects to buy $y_{im}$ outputs from provider $m$, the single-provider subscription condition is

\begin{equation}\left(P_m^{\mathrm{spot}}-\rho_m\right)\cdot\mathbb E[y_{im}]\geq\phi_m.\end{equation}

Let purchaser $i$ face task types indexed by $k$, with $T_{ik}$ expected tasks of each type. Let the net value of routing task $k$ to subscribed provider $m$ be

\begin{equation}V_{imk}=\theta_i\cdot Q_{mk}-\rho_m.\end{equation}

If $S_i$ is the purchaser's subscription portfolio, expected utility is

\begin{equation}U_i(S_i)=\sum_k T_{ik}\cdot\max\left\{0,\max_{m\in S_i}V_{imk}\right\}-\sum_{m\in S_i}\phi_m.\end{equation}

The purchaser chooses

\begin{equation}S_i^*=\arg\max_{S_i\subseteq\{1,\ldots,M\}}U_i(S_i).\end{equation}

Define the best net value currently available for task $k$ from portfolio $S$ as

\begin{equation}B_{ik}(S)=\max\left\{0,\max_{m\in S}\left(\theta_i\cdot Q_{mk}-\rho_m\right)\right\}.\end{equation}

Under deterministic quality, the incremental value of adding provider $j$ is

\begin{equation}\Omega_{ij}(S)=\sum_k T_{ik}\cdot\left[\theta_i\cdot Q_{jk}-\rho_j-B_{ik}(S)\right]_+-E_{ij},\qquad [z]_+\equiv\max\{z,0\},\end{equation}

where $Q_{jk}$ is the quality of provider $j$'s model on task $k$ (as $Q_{mk}$ above, with $m=j$) and $E_{ij}\geq0$ is a fixed evaluation cost the purchaser incurs to assess provider $j$. Note that $E_{ij}$ is a one-time, relationship-level cost, paid on adding provider $j$ to the portfolio and never again---the stock counterpart of the recurring, per-output friction $s_{imk}$ of Section~\ref{subsec:task-segmentation}. The two are distinct primitives, and the distinction carries weight: the niche-protection condition of Proposition~\ref{prop:segmented-entry} can be satisfied only by the recurring kind, precisely because a cost of the $E_{ij}$ kind is extinguished the moment a purchaser multi-homes through the mechanism of this appendix.

Purchaser $i$ adds provider $j$ when

\begin{equation}\Omega_{ij}(S)\geq\phi_j.\end{equation}

This mechanism produces multi-homing even under complete information. Provider $j$ need not be the purchaser's best provider on average; it need only be sufficiently better or cheaper for some tasks to justify its fixed fee.

If task-specific quality is uncertain but observed before routing, let $Q_{ijk}$ be random. The incremental option value becomes

\begin{equation}\Omega_{ij}(S)=\sum_k T_{ik}\cdot\mathbb E\left[\theta_i\cdot Q_{ijk}-\rho_j-B_{ik}(S)\right]_+-E_{ij}.\end{equation}

If quality is learned only by testing a provider, the testing decision requires a separate information-acquisition model rather than a fixed evaluation cost alone.

Let $n_m$ be the mass of subscribers and $Y_m$ total subscriber usage. If the wholesale compute price is $r_m$, marginal output cost is

\begin{equation}\kappa_m=b_m+a_m\cdot r_m.\end{equation}

Operating profit is

\begin{equation}\Pi_m=\phi_m\cdot n_m+\left(\rho_m-\kappa_m\right)\cdot Y_m-C_m,\end{equation}

where $C_m$ is provider $m$'s fixed cost per period.

A provider may set $\rho_m<\kappa_m$ if fixed subscription revenue covers the usage subsidy:

\begin{equation}\phi_m\cdot n_m\geq\left(\kappa_m-\rho_m\right)\cdot Y_m+C_m.\end{equation}

Multi-homing separates subscription share, usage share, and compute share. A trailing provider can have a broad subscriber base while receiving a much smaller fraction of actual inference traffic.

\section{A Symmetric Cournot Example}
\label{app:cournot}

Suppose $n$ model providers sell a homogeneous inference service with inverse demand

\begin{equation}P(Y)=A-B\cdot Y,\qquad Y=\sum_{\ell=1}^{n}y_\ell.\end{equation}

Each output requires $a$ FLOPs and has non-compute cost $b$. Given wholesale compute price $r$, marginal cost is $\kappa=b+a\cdot r$. The symmetric Cournot equilibrium is

\begin{equation}y_\ell=\frac{A-\kappa}{B\cdot(n+1)},\qquad Y=\frac{n\cdot(A-b-a\cdot r)}{B\cdot(n+1)}.\end{equation}

If aggregate capacity binds, $a\cdot Y=K$, so

\begin{equation}r^*=\frac{A-b}{a}-\frac{B\cdot(n+1)\cdot K}{a^2\cdot n}.\end{equation}

Consumer price is

\begin{equation}P^*=A-\frac{B\cdot K}{a}.\end{equation}

Conditional on full utilization of fixed capacity, consumer price depends on aggregate capacity but not directly on the number of model providers. The number of providers instead changes the division of rent. Each laboratory's markup is

\begin{equation}P^*-\kappa=\frac{B\cdot K}{a\cdot n},\end{equation}

and aggregate laboratory operating profit is

\begin{equation}\Pi^{\mathrm{labs}}=\frac{B\cdot K^2}{a^2\cdot n}.\end{equation}

As $n$ increases, laboratory oligopoly rent falls and the wholesale compute price rises, transferring more of the fixed-capacity rent to hyperscalers. At physical compute cost $r=c$, Cournot output is

\begin{equation}Y^{\mathrm{Cournot}}(c)=\frac{n\cdot(A-b-a\cdot c)}{B\cdot(n+1)}.\end{equation}

Capacity binds only when

\begin{equation}\label{eq:cournot-slack-threshold}K<\frac{a\cdot n\cdot(A-b-a\cdot c)}{B\cdot(n+1)}.\end{equation}

Above this threshold, capacity is slack, $r^*=c$, hyperscaler scarcity rents vanish, and laboratories retain ordinary Cournot markups. Appendix~\ref{app:slack-capacity}, next, takes up that slack case in full.

\section{The Slack-Capacity Regime}
\label{app:slack-capacity}

The main text, from Section~\ref{sec:fixed-capacity} onward, maintains the scarcity assumption: desired output at prices near physical production cost exceeds available capacity, the multiplier $\lambda$ is positive, and the wholesale price of compute carries a scarcity premium, $r^*>c$. The opposite regime---capacity slack, $\lambda=0$, $r^*=c$---appears in the main text only as asides attached to binding-case results: the slack branch of Section~\ref{sec:oligopoly}, the closing sentence of Proposition~\ref{prop:capacity-neutrality}'s discussion, part~(4) of Proposition~\ref{prop:free-entry}, and the scope restriction around the Bertrand analysis of Section~\ref{subsec:price-competition}. This appendix works the regime out in one place. It introduces no new primitives and no new notation; exact statements use the symmetric linear-demand system of Appendix~\ref{app:cournot}, which is the closest sibling of everything derived here. Following the standard of Section~\ref{subsec:entry-honesty}, each subsection states what is derived (within that linear system), what is a qualitative implication of derived objects, and what is left open.

\subsection{When is capacity slack?}
\label{subsec:slack-condition}

Slackness is the negation of the scarcity assumption of Section~\ref{sec:fixed-capacity}, but there is no single conduct-free version of the condition, because what must fall short of $K$ is \emph{desired} output at compute price $c$, and desired output depends on who is choosing it. The paper already contains the two operative versions:

\begin{equation}\sum_{m=1}^{M}a_m\cdot D_m(\mathbf P^0,\mathbf Q)\leq K\qquad\text{(integrated benchmark, Section~\ref{sec:fixed-capacity})},\end{equation}

\begin{equation}\sum_\ell a_\ell\cdot y_\ell^{\mathrm{olig}}(c)\leq K\qquad\text{(downstream oligopoly, Section~\ref{sec:oligopoly})}.\end{equation}

In the symmetric linear system these become explicit thresholds, derived in scattered places earlier; collecting them, write

\begin{equation}\bar K\equiv\frac{a\cdot(A-b-a\cdot c)}{B}\end{equation}

for the compute that demand at near-cost prices would absorb under price-taking conduct. Head-on Cournot among $n$ laboratories leaves capacity slack when $K\geq\frac{n}{n+1}\cdot\bar K$ (equation~\eqref{eq:cournot-slack-threshold}); the segmented configuration of Proposition~\ref{prop:segmented-entry}, whose segment monopolists individually restrict output harder, is slack already at $K\geq\frac{1}{2}\cdot\bar K$ (the threshold derived in the discussion following that proposition, and the $n=1$ case of the Cournot threshold); and the price-taking benchmark is slack only at $K\geq\bar K$.

Three consequences of this ordering are worth recording, because each is used below. First, the boundary is endogenous to conduct: for $K$ between $\frac{n}{n+1}\cdot\bar K$ and $\bar K$, capacity is slack under oligopoly but would bind under price-taking conduct---idle FLOPs coexist with unserved purchasers whose willingness to pay exceeds physical cost, because output is withheld by downstream market power rather than rationed by hardware. The observation of Section~\ref{sec:oligopoly} that additional GPU construction does not lower prices in the slack regime has its sharpest form in this window: capacity is not what is scarce there. Second, the boundary is endogenous to $n$: the Cournot threshold $\frac{n}{n+1}\cdot\bar K$ rises with $n$, so entry alone can flip a slack market back to binding, a fact the free-entry analysis below has to respect. Third, in the linear system the wholesale price reaches cost exactly at the threshold and stays there (Appendix~\ref{app:cournot}), so the two regimes join continuously rather than across a jump; there is no discontinuity for a transition to leap over.

\subsection{Cournot competition with the wholesale price pinned at cost}
\label{subsec:slack-cournot}

Setting $r=c$ in the equilibrium of Appendix~\ref{app:cournot} gives the standard unconstrained Cournot outcome. With marginal cost $\kappa=b+a\cdot c$,

\begin{equation}y_\ell=\frac{A-b-a\cdot c}{B\cdot(n+1)},\qquad Y=\frac{n\cdot(A-b-a\cdot c)}{B\cdot(n+1)},\qquad P^{\mathrm{slack}}=b+a\cdot c+\frac{A-b-a\cdot c}{n+1}.\end{equation}

The three-component price decomposition of Section~\ref{sec:oligopoly} still applies, with the middle component dead:

\begin{equation}P^{\mathrm{slack}}=\underbrace{b+a\cdot c}_{\text{physical production cost}}+\underbrace{0}_{\text{hyperscaler scarcity rent}}+\underbrace{\frac{A-b-a\cdot c}{n+1}}_{\text{model-provider markup}}.\end{equation}

The contrast with Proposition~\ref{prop:capacity-neutrality} is the spine of this appendix, and Table~\ref{tab:regime-contrast} states it side by side. Under binding capacity the price is $A-B\cdot K/a$: it depends on aggregate capacity and not on the number of laboratories, entry only reallocates rent toward hyperscalers, and adding laboratories without limit leaves the consumer price unchanged. Under slack capacity the dependence inverts completely: the price depends on $n$ and not on $K$, additional capacity has no direct price effect at all, and the $n\to\infty$ limit is marginal-cost pricing. The two policy levers of Proposition~\ref{prop:capacity-neutrality}'s discussion swap roles at the boundary: while capacity binds, building GPUs lowers prices and adding laboratories does not; once capacity is slack, adding laboratories lowers prices and building GPUs does not.

\begin{table}[htbp]
\centering
\caption{Binding versus slack capacity in the symmetric linear Cournot system (Appendix~\ref{app:cournot}), with $\bar K\equiv a\cdot(A-b-a\cdot c)/B$.}
\label{tab:regime-contrast}
\begin{tabular}{lll}
\hline
 & Binding: $K<\frac{n}{n+1}\bar K$ & Slack: $K\geq\frac{n}{n+1}\bar K$ \\
\hline
Consumer price $P^*$ & $A-B\cdot K/a$ & $b+a\cdot c+(A-b-a\cdot c)/(n+1)$ \\
$P^*$ moves with & $K$ only & $n$ only \\
Wholesale price $r^*$ & $\frac{A-b}{a}-\frac{B\cdot(n+1)\cdot K}{a^{2}\cdot n}>c$ & $c$ \\
Laboratory markup & $B\cdot K/(a\cdot n)$ & $(A-b-a\cdot c)/(n+1)$ \\
Flow profit per laboratory & $B\cdot K^{2}/(a^{2}\cdot n^{2})$ & $(A-b-a\cdot c)^{2}/(B\cdot(n+1)^{2})$ \\
Effect of entry on $P^*$ & none (rent moves upstream) & lowers it \\
Effect of added capacity on $P^*$ & lowers it & none (direct) \\
\hline
\end{tabular}
\end{table}

The regimes join continuously: at the threshold $K=\frac{n}{n+1}\cdot\bar K$, the binding-case price $A-B\cdot K/a$ equals $b+a\cdot c+(A-b-a\cdot c)/(n+1)$ exactly, so the consumer price as a function of $K$ is continuous and kinked---strictly decreasing in $K$ up to the threshold, flat beyond it. What changes at the boundary is which derivative is nonzero, not the level. All of this is derived rather than asserted, but derived for the symmetric homogeneous-good linear system: with differentiated products or heterogeneous laboratories the slack-regime price still depends on $n$ in the ordinary oligopoly way, but the formulas above are no longer exact.

\subsection{Free entry under slack capacity}
\label{subsec:slack-free-entry}

Part~(4) of Proposition~\ref{prop:free-entry} stated the slack free-entry count in one line; here is the derivation. With capacity slack, an active laboratory earns markup times output from the previous subsection,

\begin{equation}\pi^{\mathrm{slack}}(n)=\frac{(A-b-a\cdot c)^{2}}{B\cdot(n+1)^{2}}.\end{equation}

The entry condition $\pi^{\mathrm{slack}}(n)/(\rho+\Lambda)\geq\mathcal F$ rearranges to

\begin{equation}n+1\leq\frac{A-b-a\cdot c}{\sqrt{B\cdot(\rho+\Lambda)\cdot\mathcal F}},\end{equation}

so the free-entry count is $n^{*}_{\mathrm{slack}}+1=\left\lfloor(A-b-a\cdot c)/\sqrt{B\cdot(\rho+\Lambda)\cdot\mathcal F}\right\rfloor$, the expression quoted in Proposition~\ref{prop:free-entry}(4). Substituting the uniform-$\theta$ microfoundation $A=\bar\theta\cdot Q$ and $B=\bar\theta\cdot Q/N$ gives

\begin{equation}n^{*}_{\mathrm{slack}}+1=\left\lfloor\left(\bar\theta\cdot Q-b-a\cdot c\right)\cdot\sqrt{\frac{N}{\bar\theta\cdot Q\cdot(\rho+\Lambda)\cdot\mathcal F}}\right\rfloor,\end{equation}

which rises with the purchaser mass proportionally to $\sqrt N$, as asserted there. The economics is the classic sunk-cost logic of \citet{sutton1991} operating unobstructed: under slack capacity the aggregate laboratory profit pool $n\cdot\pi^{\mathrm{slack}}(n)$ scales linearly with $N$ (because $B\propto1/N$), so a market twice as large carries $\sqrt2$ times as many laboratories past a given fixed-cost bar. Under binding capacity the same calculation runs backward---the pool $B\cdot K^{2}/(a^{2}\cdot n)$ \emph{shrinks} with $N$---which is the reversal emphasized after Proposition~\ref{prop:free-entry}: market growth funds entry only on the slack side of the boundary.

The two counts combine more tightly than a case split suggests. At any $(n,K)$ the operative flow profit is the smaller of the two expressions---when capacity binds at $n$ laboratories the binding-case profit lies below the slack-case one, and conversely---so flow profit is $\min\{B\cdot K^{2}/(a^{2}\cdot n^{2}),\,(A-b-a\cdot c)^{2}/(B\cdot(n+1)^{2})\}$, continuous in both arguments and strictly decreasing in $n$ in both branches, and the free-entry count at any $K$ is simply

\begin{equation}n^{*}=\min\left\{n^{*}_{\mathrm{bind}},\,n^{*}_{\mathrm{slack}}\right\},\end{equation}

with $n^{*}_{\mathrm{bind}}$ the binding-capacity count of Proposition~\ref{prop:free-entry}(1). This resolves a consistency question the one-line aside left open: the regime depends on $n$ (Section~\ref{subsec:slack-condition}) while $n$ depends on the regime, but because profit is continuous across the boundary and monotone in $n$, the fixed point is the minimum of the two counts, with no cycling and no gap. One comparative static of the combined object is striking enough to record, with its scope stated plainly. Holding $K$ fixed and varying market size, $n^{*}_{\mathrm{slack}}$ rises with $N$ while $n^{*}_{\mathrm{bind}}$ falls with it, so the equilibrium count is single-peaked in $N$ and is maximized where the two bounds cross---at market sizes that put the industry near the binding/slack boundary itself. Small markets support few laboratories because the profit pool is small; large markets support few because the pool drains upstream to hyperscalers; the count peaks in between. We flag this as a derived property of the symmetric linear example, not a general theorem, and it inherits every caveat of Section~\ref{subsec:entry-honesty}, in particular the unmodeled escalation of entry quality.

\subsection{Why investment should not aim past the boundary}
\label{subsec:slack-investment}

The hyperscaler capital-expenditure condition of Section~\ref{sec:model-race} prices future capacity against expected future scarcity rents $\lambda_{h,t+s}$. In the linear system those rents have an explicit form: for $K$ below the $n$-laboratory Cournot threshold,

\begin{equation}\lambda(K)=r^{*}(K)-c=\frac{B\cdot(n+1)}{a^{2}\cdot n}\cdot\left(\frac{n}{n+1}\cdot\bar K-K\right),\end{equation}

linear and strictly decreasing in $K$, reaching zero exactly at the threshold and remaining zero beyond it. The rent that rewards capacity investment does not fall off a cliff at the boundary; it shrinks continuously to nothing as capacity approaches it.

Two implications follow, one derived and one qualitative, and the difference in status matters. The derived implication concerns a deterministic stationary environment. With depreciation $\delta_K>0$, maintaining any capacity level requires perpetual positive investment, and the investment condition of Section~\ref{sec:model-race} equates marginal capital cost to the discounted stream of rents a delivered FLOP of capacity will earn---roughly $\lambda_{ss}\cdot\beta^{T_{\mathrm{build}}}\cdot\Phi_h'/\left(1-\beta\cdot(1-\delta_K)\right)=C_h'(H_{ss})$. A steady state in the slack region would set the left side to zero against a strictly positive right side, a contradiction; so any deterministic stationary equilibrium with costly capacity sits strictly inside the binding region, with the steady-state scarcity rent equal to the amortized cost of capacity. This is the peak-load-pricing logic of \citet{boiteux1960} and \citet{crewKleindorfer1976}, cited in Section~\ref{sec:oligopoly}, recovered from the model's own investment condition: planned, durable slack is inconsistent with hyperscalers choosing to replace depreciating hardware, because at $r^{*}=c$ a FLOP of capacity recovers its operating cost and none of its capital cost.

The qualitative implication is about the world outside that stationary benchmark, and it should be stated with care rather than oversold. Nothing above says slack episodes will not happen. Capacity is chosen under uncertainty and arrives after the lag $T_{\mathrm{build}}$, so hardware ordered when expected rents were high lands in whatever world exists on delivery; a disappointing release cycle, a pause in demand growth, or a burst of inference-efficiency gains (a falling $a$) can each push a market whose capacity was rationally planned into realized slack. The investment condition disciplines the expectation of $\lambda$, not its realization: transient slack is the bad tail of an ex ante sensible investment program, and its occurrence is not evidence of a mistake. What the model does say is that slack should be self-correcting on the investment margin---capacity growth decelerates as the boundary nears, because the marginal rent backing it shrinks---and on the demand margin whenever demand or training appetite grows. Three channels could defeat even this weaker claim, and all sit outside the model: strategic investment races among hyperscalers, whose joint capacity choices are not solved here (the motivation for citing \citet{besankoDoraszelski2004} in Section~\ref{sec:literature-review}); captive internal demand of the integrated suppliers of Section~\ref{sec:vertical-integration}, which decouples investment from the market rent $\lambda$; and any non-rent motive for building, from preemption to option value on future demand regimes. The honest summary: that $\lambda$ vanishes continuously at the boundary is derived; that overshoot is unlikely to be \emph{planned} is a qualitative implication of that derivation; that overshoot will not \emph{happen} is not a claim the model supports.

\subsection{Niche protection without scarcity rents}
\label{subsec:slack-niches}

Proposition~\ref{prop:segmented-entry} priced niche protection under binding capacity: the recurring friction had to absorb the full segment markup, $s_{\min}=B\cdot K/a$, and part of what made that markup no larger was that a rival's alternative to invading was selling output elsewhere at the scarcity price $r^{*}>c$. Under slack capacity the same three-line derivation runs with $r=c$. A segment monopolist facing marginal cost $b+a\cdot c$ sets $A-2\cdot S\cdot B\cdot y_k=b+a\cdot c$, so

\begin{equation}y_k=\frac{A-b-a\cdot c}{2\cdot S\cdot B},\qquad P_k^{\mathrm{slack}}=b+a\cdot c+\frac{A-b-a\cdot c}{2},\qquad \pi^{\mathrm{seg,slack}}(S)=\frac{(A-b-a\cdot c)^{2}}{4\cdot S\cdot B},\end{equation}

with the markup $(A-b-a\cdot c)/2$ now the unconstrained monopoly markup, independent of both $S$ and $K$. The protection condition is the same participation logic as in the binding case: a rival's first unit in segment $k$ is unprofitable when $(P_k^{\mathrm{slack}}-s-b)/a\leq r^{*}=c$, which rearranges to

\begin{equation}s\;\geq\;\frac{A-b-a\cdot c}{2},\end{equation}

and the contested-segment equilibrium confirms the threshold from the other side exactly as before: with $r$ pinned at $c$ there is no capacity-clearing feedback, the entrant's output is $y_E=(A-b-a\cdot c-2\cdot s)/(3\cdot S\cdot B)$, and its flow profit $(A-b-a\cdot c-2\cdot s)^{2}/(9\cdot S\cdot B)$ is positive below the threshold and vanishes exactly at it.

Three things change relative to the binding case, and one does not. First, the required friction is larger. On the binding domain $s_{\min}=B\cdot K/a$ rises with $K$ and reaches $(A-b-a\cdot c)/2$ exactly at the segmented threshold $K=\frac12\cdot\bar K$, where it plateaus: the required friction is continuous in $K$, increasing while capacity binds, and maximal under slack. The economic reading is that compute scarcity itself was doing part of the protective work---a high $r^{*}$ made a rival's first invading FLOP expensive---and as the scarcity rent erodes, the friction must absorb the whole of a now-larger margin. Capacity growth therefore erodes niche protection continuously, well before the regime flips. Second, segmentation stops being price-neutral. Under binding capacity, part~(2) of Proposition~\ref{prop:segmented-entry} showed that segmentation leaves the consumer price unchanged and only redistributes rent; under slack, the segmented price $b+a\cdot c+(A-b-a\cdot c)/2$ strictly exceeds the head-on Cournot price $b+a\cdot c+(A-b-a\cdot c)/(S+1)$ for any $S\geq2$. Market division into protected monopolies now raises prices and restricts output---the ordinary welfare cost of market division, which capacity-pinned output previously suppressed, in line with the form-switching observation of Section~\ref{sec:welfare}. Third, the viability arithmetic survives in parallel form: free entry supports up to $\bar S_{\mathrm{slack}}=\left\lfloor(A-b-a\cdot c)^{2}/(4\cdot B\cdot(\rho+\Lambda)\cdot\mathcal F)\right\rfloor$ segment monopolists, which is (up to integer parts) the square of the head-on slack count divided by four, so fragmentation into defensible niches still multiplies the number of laboratories a given profit pool can carry.

What does not change is the honest status of the whole construction, and here the slack case is if anything weaker. The friction magnitude now required---half the choke-price margin, regardless of capacity---is even more demanding than the binding-case threshold, and the paper offers no account of any real recurring friction reaching it. Without such a friction, quality parity plus slack capacity is the textbook homogeneous-good case: Bertrand undercutting competes markups to zero and no niche survives. So the direct answer to whether compute-scarcity-based protection disappears is yes, and quantifiably: the protective wedge shrinks one-for-one with the scarcity rent, and at $r^{*}=c$ the friction is the entire defense. Whether ordinary horizontal product differentiation---taste-based rather than friction-based---would step into the role that scarcity vacates is a standard question of differentiated-products oligopoly \citep{singhVives1984}, but it is not answerable within this paper's primitives, which contain vertical quality, a recurring friction wedge, and nothing else horizontal; we leave it genuinely open. And the assignment-selection gap discussed after Proposition~\ref{prop:segmented-entry} carries over verbatim: nothing here explains why $S$ laboratories each pay $\mathcal F$ rather than one laboratory spanning every segment.

\subsection{Observable markers of the slack regime}
\label{subsec:slack-empirics}

Which regime a market is in is, in principle, observable, and the two regimes disagree about enough that the distinction is testable rather than rhetorical. The direct marker is the wholesale margin $r^{*}-c$: rental prices for inference-capable compute at physical operating cost---contributing nothing to capital recovery---signal slack, as does persistent measured utilization below capacity. The indirect markers are the comparative statics of Table~\ref{tab:regime-contrast}: in the scarcity regime the consumer price of a quality tier tracks aggregate GPU supply and ignores entry (empirical prediction~1 of Section~\ref{sec:empirical}), while in the slack regime it tracks entry and near-frontier releases and ignores capacity growth. Prediction~5 of Section~\ref{sec:empirical} states the resulting regime-switch test. The transition itself has identifiable drivers on both sides: demand growth and rising training appetite push toward binding; capacity construction and inference-efficiency gains (falling $a$) push toward slack; and---less obviously, from Section~\ref{subsec:slack-condition}---entry pushes toward binding, because more laboratories jointly restrict output less. A market can therefore cross the boundary in either direction without any change in the physical capacity stock.

Two things this appendix deliberately does not claim. The single-peaked laboratory-count property of Section~\ref{subsec:slack-free-entry} is a feature of one symmetric linear example, and we do not state it as an empirical prediction. And nothing here forecasts which regime the industry will occupy at any horizon: the investment analysis implies the market should hover near the binding side of the boundary in expectation, with transient slack episodes, but that is a statement about incentives around the boundary, not a calibration of where the boundary currently sits.

\section{Task Heterogeneity, Unit Invariance, and Multiple Resource Constraints}
\label{app:extensions}

\subsection{Task-specific productivity}

If model performance differs across tasks, let $Q_{m,k}$ denote model $m$'s economically relevant quality on task $k$, and define

\begin{equation}q_{m,k}=\frac{Q_{m,k}}{a_{m,k}}.\end{equation}

Purchaser $i$, whose task is $k_i$, receives utility

\begin{equation}u_{im}=\theta_i\cdot Q_{m,k_i}-P_{m,k_i}.\end{equation}

A model may trail in general capability while generating the highest marginal net revenue per FLOP on a particular class of tasks. This provides a natural basis for coexistence without requiring a single scalar ranking of all models.

\subsection{Output-unit invariance}

A pure change in the definition of a standardized output unit should not alter the economic allocation. Consider a rescaling by $\gamma>0$:

\begin{equation}a_m'=\gamma\cdot a_m,\qquad Q_m'=\gamma\cdot Q_m,\qquad P_m'=\gamma\cdot P_m,\qquad b_m'=\gamma\cdot b_m,\qquad y_m'=\frac{y_m}{\gamma}.\end{equation}

Then aggregate compute, revenue, cost, and quality are unchanged:

\begin{equation}a_m'\cdot y_m'=a_m\cdot y_m,\qquad P_m'\cdot y_m'=P_m\cdot y_m,\qquad Q_m'\cdot y_m'=Q_m\cdot y_m.\end{equation}

The normalized quantities are also unchanged:

\begin{equation}q_m'=q_m,\qquad p_m'=p_m,\qquad \frac{b_m'}{a_m'}=\frac{b_m}{a_m}.\end{equation}

Thus merely redefining the output unit does not create an efficiency advantage. A genuine reduction in inference requirements matters when it preserves more than a proportional share of purchaser value or lowers cost per unit of economically relevant quality.

\subsection{Multiple constrained resources}

FLOPs may not be the only binding resource. Memory capacity, memory bandwidth, latency, networking, power, or accelerator-specific availability may constrain inference independently.

Let $a_m^r$ be the amount of resource $r$ required per output, let $K^r$ be aggregate capacity of that resource, and let $c_r$ be its physical marginal cost. The allocation problem becomes

\begin{equation}\max_{\mathbf y\geq0}\;R(\mathbf y;\mathbf Q)-\sum_m b_m\cdot y_m-\sum_r c_r\cdot\sum_m a_m^r\cdot y_m\quad\text{subject to}\quad\sum_m a_m^r\cdot y_m\leq K^r\;\text{for each }r.\end{equation}

If $\lambda_r$ is the shadow price of resource $r$, an active model satisfies

\begin{equation}MR_m-b_m=\sum_r\left(c_r+\lambda_r\right)\cdot a_m^r.\end{equation}

The general allocation rule is therefore to compare the marginal revenue created by additional inference effort with the shadow-priced vector of constrained resources, rather than with FLOPs alone.

\subsection{Nonseparable downstream demand}

If one laboratory's output directly affects the marginal value of another laboratory's output, write laboratory $\ell$'s marginal value as

\begin{equation}v_\ell(x_\ell,\mathbf x_{-\ell}).\end{equation}

A candidate allocation in which leader $L$ receives all capacity is optimal only if

\begin{equation}v_L(K,\mathbf 0_{-L})\geq v_j\left(0;x_L=K,\mathbf 0_{-(L,j)}\right)\end{equation}

for every follower $j$. This is the nonseparable counterpart to the simpler main-text condition $v_L(K)\geq\max_j v_j(0)$.

\section{Two-Model Demand System: Derivations}
\label{app:two-model}

This appendix collects the calculations behind Lemma~\ref{lem:two-model} and Proposition~\ref{prop:breakthrough}. The segmentation, demand system, inverse demands, and revenue function are stated in Section~\ref{subsec:two-model-demand}; they are taken as given here.

\subsection{Concavity of the revenue function}

Differentiating $R(\mathbf y;\mathbf Q)=\bar\theta(Q_L y_L+Q_F y_F)-(\bar\theta/N)(Q_L y_L^{2}+2Q_F y_L y_F+Q_F y_F^{2})$ twice gives the constant Hessian

\begin{equation}\nabla^{2}R=-\frac{2\bar\theta}{N}\begin{pmatrix}Q_L & Q_F\\ Q_F & Q_F\end{pmatrix}.\end{equation}

Its leading principal minor is $-2\bar\theta Q_L/N<0$, and its determinant is $(2\bar\theta/N)^{2}\,Q_F(Q_L-Q_F)>0$ because $Q_L>Q_F$. The Hessian is therefore negative definite, so $R$ is strictly concave wherever the three-tier segmentation holds. This is part~(1) of Lemma~\ref{lem:two-model}.

\subsection{Productivity and marginal net revenue per FLOP}

The marginal revenues are

\begin{equation}MR_L=\frac{\partial R}{\partial y_L}=\bar\theta\,Q_L-\frac{2\bar\theta}{N}\left(Q_L y_L+Q_F y_F\right),\qquad MR_F=\frac{\partial R}{\partial y_F}=\bar\theta\,Q_F\left[1-\frac{2(y_L+y_F)}{N}\right].\end{equation}

Writing $Q_m=a_m q_m$ and $v_m=(MR_m-b_m)/a_m$, and holding output quantities fixed,

\begin{equation}\frac{\partial v_L}{\partial q_L}=\bar\theta\left(1-\frac{2y_L}{N}\right),\qquad \frac{\partial v_F}{\partial q_F}=\bar\theta\left(1-\frac{2(y_L+y_F)}{N}\right).\end{equation}

Both are strictly positive when the served mass $y_L+y_F$ is below $N/2$, i.e., when capacity is scarce enough that only the higher-$\theta$ purchasers are served. This is part~(2) of Lemma~\ref{lem:two-model}.

\subsection{Bertrand equilibrium and comparative statics (slack capacity): Proposition~\ref{prop:breakthrough}, part 1}
\label{subsec:app-bertrand}

This subsection proves part~1 of Proposition~\ref{prop:breakthrough}. It is a slack-capacity local analysis: laboratories set prices and each obtains, at that price, whatever compute the resulting quantity requires at the constant compute-inclusive marginal cost $\kappa$. Nothing here rations output below what unconstrained price competition would choose; the capacity-constrained case is treated separately in the next subsection, which is where part~2 of the proposition, and the paper's maintained scarcity regime, actually live (Section~\ref{sec:limitations} states this scope restriction explicitly).

Let laboratories set prices with constant compute-inclusive marginal costs $\kappa_L,\kappa_F$. The leader's residual demand slope,

\begin{equation}\frac{\partial D_L}{\partial P_L}=-\frac{N}{(Q_L-Q_F)\,\bar\theta},\end{equation}

is a primitive of the demand system, requiring no equilibrium argument: its magnitude is strictly decreasing in the quality gap $Q_L-Q_F$, so a wider gap makes the leader's residual demand less price-sensitive, and closer substitutes (a smaller gap) make it more sensitive. The interior first-order condition gives the leader's Lerner markup $\mu_L\equiv P_L-\kappa_L=D_L\big/\lvert\partial D_L/\partial P_L\rvert$.

Take the scale normalization $N=\bar\theta=1$ (a normalization of units, not of cost) and let $\Delta\equiv Q_L-Q_F>0$. Solving both first-order conditions with common compute-inclusive cost $\kappa_L=\kappa_F=\kappa\geq0$ gives

\begin{equation}P_L=\frac{Q_L(2\Delta+3\kappa)}{3Q_L+\Delta},\qquad P_F=\frac{(Q_L-\Delta)P_L}{2Q_L}+\frac{\kappa}{2},\end{equation}

so that the leader's markup is

\begin{equation}\mu_L\equiv P_L-\kappa=\frac{\Delta\,(2Q_L-\kappa)}{3Q_L+\Delta}.\end{equation}

At $\kappa=0$ this reduces to the zero-cost case worked in earlier drafts of this appendix, $P_L=2Q_L\Delta/(4Q_L-\Delta)$ and $\mu_L=P_L$; the derivation above keeps $\kappa$ general precisely because the zero-cost case alone cannot support the elasticity claim below. Direct differentiation of $\mu_L$ with respect to $\Delta$, holding $Q_L$ and $\kappa$ fixed, gives

\begin{equation}\frac{\partial\mu_L}{\partial\Delta}=\frac{(2Q_L-\kappa)\cdot 3Q_L}{(3Q_L+\Delta)^{2}},\end{equation}

which is strictly positive whenever $2Q_L>\kappa$---equivalently, whenever the equilibrium is interior with $\mu_L>0$, the case the proposition restricts to. So $\mu_L$ is strictly increasing in the quality gap, and $\mu_L\to0$ as $\Delta\to0$. This is the first half of part~1.

For the elasticity, the Lerner identity gives $\epsilon_{LL}=P_L/\mu_L$. Substituting the expressions above,

\begin{equation}\epsilon_{LL}=\frac{Q_L(2\Delta+3\kappa)}{\Delta(2Q_L-\kappa)}=\frac{Q_L}{2Q_L-\kappa}\left(2+\frac{3\kappa}{\Delta}\right).\end{equation}

Differentiating with respect to $\Delta$ holding $Q_L,\kappa$ fixed,

\begin{equation}\frac{\partial\epsilon_{LL}}{\partial\Delta}=-\frac{3\kappa\,Q_L}{(2Q_L-\kappa)\,\Delta^{2}}.\end{equation}

This is strictly negative whenever $\kappa>0$: with any positive compute-inclusive cost, the leader's equilibrium own-price elasticity is strictly decreasing in the quality gap, which is exactly the elasticity claim of part~1, now shown by differentiation rather than asserted. At $\kappa=0$ the derivative is identically zero and $\epsilon_{LL}\equiv1$ for every $\Delta$---the case worked in earlier drafts of this appendix is the degenerate boundary of this result, not a demonstration of it, which is why the elasticity claim in the proposition is stated for $\kappa>0$. These are the standard comparative statics of the uncovered-market vertical-differentiation duopoly \citep{shakedSutton1982,singhVives1984}, and they are exactly the first arrow chain of Section~\ref{subsec:continuing-improvement}.

\subsection{Capacity-constrained Cournot and the compute-rent comparative static: Proposition~\ref{prop:breakthrough}, part 2}
\label{subsec:app-cournot-breakthrough}

Part~1 above is silent about $\widehat v_L-\widehat v_F$ because unconstrained Bertrand competition has no compute endowment in it: $\kappa$ is a fixed parameter, not the outcome of a binding constraint, so there is no $x_\ell$ for $\partial\pi^*_\ell/\partial x_\ell$ to be a derivative with respect to. This subsection builds the compute-rent claim in a model where $x_\ell$ actually appears, using the same two-model demand system.

\citet{krepsScheinkman1983} show that when firms first commit capacity and then compete in prices subject to that capacity (with efficient rationing), the resulting outcome coincides with the Cournot quantity outcome at those capacities. We use this to justify treating the capacity-constrained regime as a quantity game directly: laboratory $\ell$ precommits compute $x_\ell$, deploys it as output $y_\ell=x_\ell/a_\ell$, and receives the price read off the inverse demand system of Section~\ref{subsec:two-model-demand} at $(y_L,y_F)$. With $a_L=a_F=1$ and $N=\bar\theta=1$ (again a units normalization), $y_\ell=x_\ell$ and

\begin{equation}P_L(y_L,y_F)=Q_L-(Q_L y_L+Q_F y_F),\qquad P_F(y_L,y_F)=Q_F(1-y_L-y_F).\end{equation}

Laboratory $\ell$'s realized operating profit, net of the wholesale cost $r_\ell$ per FLOP folded into $\kappa_\ell\equiv b_\ell+r_\ell$ (using $a_\ell=1$), is $\pi_\ell(x_\ell,x_{-\ell})=\left[P_\ell(y_L,y_F)-\kappa_\ell\right]y_\ell$ with $y_\ell=x_\ell$. This is a genuine function of $x_\ell$, so the private marginal value of compute defined in equation~\eqref{eq:vhat-def} is well posed here:

\begin{equation}\widehat v_\ell(x_\ell,x_{-\ell})=\frac{\partial\pi_\ell}{\partial x_\ell}=P_\ell(y_L,y_F)-\kappa_\ell+y_\ell\cdot\frac{\partial P_\ell}{\partial y_\ell},\qquad \frac{\partial P_L}{\partial y_L}=-Q_L,\quad\frac{\partial P_F}{\partial y_F}=-Q_F.\end{equation}

This formula holds at any allocation $(x_L,x_F)$, not only at a Nash equilibrium of the capacity game; we evaluate it at the winner-take-all corner $x_L=K$, $x_F=0$, which is the case the paper's own definition of $r^*\simeq\max\{c,\widehat v_F\}$ (Section~\ref{sec:oligopoly}) already presumes, since that expression compares the leader's marginal value to the \emph{best losing} laboratory's value at zero. At this corner,

\begin{equation}P_L(K,0)=Q_L(1-K),\qquad P_F(K,0)=Q_F(1-K),\end{equation}

so that, with common compute-inclusive cost $\kappa_L=\kappa_F=\kappa$,

\begin{equation}\widehat v_L(K,0)=Q_L(1-2K)-\kappa,\qquad \widehat v_F(0;K)=Q_F(1-K)-\kappa,\end{equation}

where the strategic term $y_F\cdot\partial P_F/\partial y_F$ drops out of $\widehat v_F$ because $y_F=0$: the follower's value at the margin is just its price net of cost, with no cannibalization term, matching the inactive-laboratory condition $v_\ell(0)\leq r^*$ used throughout Sections~\ref{sec:concentration}--\ref{sec:trailing-labs}. Subtracting, and writing $Q_F=Q_L-\Delta$,

\begin{equation}\widehat v_L-\widehat v_F=Q_L(1-2K)-Q_F(1-K)=-K\,Q_L+\Delta\,(1-K).\end{equation}

Differentiating with respect to $\Delta$ holding $Q_L$ and $\kappa$ fixed,

\begin{equation}\frac{\partial(\widehat v_L-\widehat v_F)}{\partial\Delta}=1-K,\end{equation}

strictly positive whenever $K<1$. Because $y_L=K$ cannot exceed the total purchaser mass $N=1$ under this normalization, $K<1$ is precisely the condition that capacity is insufficient to serve the entire addressable market at these prices---the scarcity assumption maintained throughout the paper (Section~\ref{sec:fixed-capacity}), not an extra restriction imposed for this result. This proves part~2 of Proposition~\ref{prop:breakthrough}: the differential private value of compute is strictly increasing in the quality gap, evaluated at the winner-take-all corner, for any common compute-inclusive cost level and any capacity level consistent with the paper's own scarcity regime. Note also that $\kappa$ cancels entirely from $\widehat v_L-\widehat v_F$ when $\kappa_L=\kappa_F$: the comparative static does not depend on the \emph{level} of the compute-inclusive cost, only on costs being equal and held fixed as $\Delta$ varies, which is a much weaker requirement than the zero-cost normalization used in part~1.

The formula also disciplines the winner-take-all corner itself: at $\Delta=0$, $\widehat v_L-\widehat v_F=-KQ_L<0$, correctly signaling that a laboratory should \emph{not} receive the entire capacity absent a quality gap; the differential only turns positive once $\Delta$ is large enough, consistent with Section~\ref{sec:concentration}'s boxed winner-take-all condition.

On the ceteris-paribus qualification flagged in the proposition: $\kappa$ is held fixed here as $\Delta$ varies, while the paper's own dynamic story (Section~\ref{sec:model-race}, empirical prediction~4) has a follower release raise $\widehat v_F$ and hence $r^*$, which would raise $\kappa$ through $r_\ell=r^*$. A full treatment would let $\kappa$ respond to $\Delta$ along the equilibrium path and solve the resulting fixed point---the same kind of fixed point the paper declines to solve for the endogenous hazard $\Lambda_L$ in Section~\ref{sec:model-race}. We do not solve it here either. The partial-equilibrium version proved above remains informative for two reasons. First, it isolates the direct, mechanical channel---a wider gap mechanically raises the leader's realized margin on its full capacity allocation and lowers the follower's foregone margin on its first unit---separately from the indirect channel running through the wholesale price, so a researcher can sign the two channels separately rather than only their sum. Second, over the short interval immediately following a release (the window the empirical predictions of Section~\ref{sec:empirical} target), $\kappa$ adjusts only as fast as the wholesale compute market re-prices, which is plausibly slower than the discrete jump in quality itself, so the fixed-$\kappa$ comparative static approximates the impact effect even though it does not describe the eventual new steady state.

\bibliographystyle{apalike}
\bibliography{paper}

\end{document}